\documentclass[10pt,a4paper]{article}
\usepackage[margin=25mm]{geometry}
\input{common/macros.tex}
\usepackage{microtype,mathtools,array}
\usepackage[hidelinks]{hyperref}

\numberwithin{equation}{section}
\setcounter{tocdepth}{2}

\theoremstyle{plain}
\newtheorem{theorem}{Theorem}[section]
\newtheorem{proposition}[theorem]{Proposition}
\newtheorem{corollary}[theorem]{Corollary}

\theoremstyle{definition}
\newtheorem{definition}[theorem]{Definition}
\newtheorem{example}[theorem]{Example}
\newtheorem{remark}[theorem]{Remark}

\hypersetup{
  pdftitle={Building Activity Representations with Pascal-Coded Signatures: A Mathematical and Experimental Companion},
  pdfauthor={Erlend Aune; Fabian N. Harang}
}

\title{Building Activity Representations with Pascal-Coded Signatures:\\[.3em]
\large A Mathematical and Experimental Companion}
\author{Erlend Aune$^{1}$ \quad Fabian N. Harang$^{2}$\\[.4em]
\small $^{1}$Cognite; Norwegian University of Science and Technology (NTNU), Trondheim\\
\small $^{2}$BI Norwegian Business School;\\
\small Centre for Sustainable, Risk-Averse and Ethical Artificial Intelligence\\
\small (SURE-AI), Simula Research Laboratory}
\date{November 2026}

\begin{document}
\maketitle

\phantomsection
\section*{Introduction}
\addcontentsline{toc}{section}{Introduction}

Sequential observations from human activity combine three qualitatively different
phenomena: static postures that persist across an observation window, transient
movements distinguished by the ordering and timing of displacements, and periodic
motions whose energy concentrates at characteristic frequencies. When such
sequences are mapped to features through path signatures and classified by a
linear readout, two design questions govern what the classifier can see. First,
what physical or statistical quantity should each discrete path increment
represent? Second, how can positional spacing along a sequence---whether across
time steps or across frequency bins---be exposed at signature depth two without
expanding ordinary signatures to high truncation depths?

This document is a standalone mathematical and experimental companion to
\emph{Building Activity Representations with Pascal-Coded
Signatures}~\cite{mainpaper}, fleshing out the mathematical details, worked
examples, feature specifications, diagnostic controls, and reproduction guides
behind the main paper. The exposition is self-contained: every path
construction, signature coordinate, algebraic identity, and experimental
protocol is defined from first principles, assuming only familiarity with
multivariable calculus, linear algebra, and the basic definition of
iterated-integral path signatures.

\paragraph{Relation to the main paper.}
Readers consulting this companion alongside~\cite{mainpaper} can map the two
texts directly as follows:
\begin{itemize}
\item \emph{Increment coding and depth-two separation} (\cite[Sections~1--2,
Fig.~1, Eqs.~(1)--(4)]{mainpaper}): Section~\ref{app:construction} introduces
depth-two tensor and logsignature coordinates, defines the Pascal increment
lift, works through the two-rise hand-movement example in detail, and formulates
observation conventions for movement, posture, and frequency spectra
(\cite[Section~3, Eq.~(5)]{mainpaper}).
\item \emph{Polynomial bases, shift algebra, and Volterra structure}
(\cite[Section~2, concluding paragraph; Section~4, opening paragraph]{mainpaper}):
Section~\ref{app:offsets} interprets the lift as a linear compression of exact
slot labels, establishes equivalence with monomial bases and the canonical
positive-diagonal QR orthogonal basis, derives the Vandermonde shift group and
twisted Chen concatenation law, and identifies the lift at all depths as a
position-only Volterra signature.
\item \emph{Comparison with higher-depth ordinary signatures}
(\cite[Section~2 and Section~3, concluding paragraph]{mainpaper}):
Section~\ref{app:depth-and-ridge} derives the piecewise-linear interpolation
correction between depth-two coded areas and depth-three clocked signatures,
exhibits an explicit depth-three collision counterexample, proves the exact
depth-\((2P-1)\) scalar staircase projection theorem, and analyzes how ridge
regularization interacts with fixed-endpoint coordinate transformations.
\item \emph{Smartphone activity recognition}
(\cite[Sections~4--5, Table~1]{mainpaper}): Section~\ref{app:har} gives the
complete specification of the \emph{compact368} representation on the official
subject-disjoint UCI HAR split, evaluates standalone blocks and single-block
deletions, and examines conventional lead--lag and spectral-staircase
replacements, closing with a numerical verification of the depth-five staircase
projection.
\item \emph{Landmark gesture trajectories}
(\cite[Section~5, Table~2]{mainpaper}): Section~\ref{app:natops} specifies the
anchored displacement path on NATOPS, organizes the nine evaluated feature maps
as a nested filtration across full and local channel-interaction graphs, and
reports cross-validation behavior, selected penalties, and paired sequence
comparisons.
\item \emph{Additional transfer and computational studies}:
Section~\ref{app:additional} extends the empirical analysis beyond~\cite{mainpaper}
with two focused studies: cross-participant transfer on the 15-sensor DSADS
benchmark and an analysis of the subject-mixed random-window UCI HAR protocol
and single-thread CPU pipeline runtime relative to SpikeDCL. The end matter
provides a unified index mapping paper and companion claims to executable
reproduction scripts and archived evidence.
\end{itemize}

\paragraph{Mathematical status and empirical scope.}
The two parts of this companion establish different kinds of conclusions, and
it is useful to state that boundary once at the outset. Part~I
(Sections~\ref{app:construction}--\ref{app:depth-and-ridge}) proves exact
algebraic and analytical statements about discrete increment lifts, polynomial
basis changes, concatenated segments, and iterated integrals. These results
characterize what chosen coordinates represent and which linear spans coincide;
they do not imply universal reductions in feature count or runtime, since
lifting \(C\) input channels by \(P\) polynomial copies produces \(PC\) channels
and quadratically many pairwise areas.
Part~I combines established signature and polynomial identities with explicit
derivations for the increment lifts studied here.

Part~II (Sections~\ref{app:har}--\ref{app:additional}) studies how these
representation choices interact with linear ridge classifiers on public
activity and gesture datasets. Throughout every experiment, all data-dependent
quantities---axis standardization, triad RMS calibration, constant-column
variance masks, final feature scaling, and cross-validated ridge penalties---are
fitted strictly on training partitions without access to held-out test windows.
At the same time, because these public benchmarks informed the iterative design
of the feature families, all empirical results should be read as exploratory
studies of representation geometry rather than pre-registered confirmatory
trials.

\paragraph{Notation and standing conventions.}
Unless stated otherwise, a discrete sequence of length \(R \geq 2\) is indexed
from zero by \(r \in \{0, \ldots, R-1\}\), with normalized grid points
\(u_r := r/(R-1) \in [0, 1]\). For scalar data, \(z=(z_r)_{r=0}^{R-1}\)
lists real-valued path increments, rather than accumulated positions.
Vectors in \(\R^C\) or \(\R^d\) carry the standard
Euclidean inner product \(\langle u, v\rangle = \sum_i u_i v_i\) and norm
\(\|u\| = \sqrt{\langle u, u\rangle}\); for a matrix \(M\), \(\|M\|_F\) denotes
the Frobenius norm \(\sqrt{\sum_{i,j} M_{ij}^2}\). Letters in a signature word
\(i_1 \cdots i_k\) are written in chronological integration order from left to
right, so \(S^{ij}(y) = \int_{0 < t_1 < t_2 < R} dy_{t_1}^i \, dy_{t_2}^j\).
For any finite sequence \(z_0, \ldots, z_{N-1} \in \R\), sample means and
standard deviations use the population normalization
\(\bar z = N^{-1}\sum_{r=0}^{N-1} z_r\) and
\(\sigma_z = \sqrt{N^{-1}\sum_{r=0}^{N-1}(z_r - \bar z)^2}\).

\clearpage
\tableofcontents
\clearpage

\part{Mathematical Foundations}
\section{Increment design and the Pascal lift}\label{app:construction}

\subsection{A motivating distinction between two movement schedules}
The net displacement of a trajectory tells how far a body part moved, but not
when the movement occurred or how separate movements were spaced across an
interval. Consider the two scalar hand-height trajectories on \([0,1]\) shown in
Figure~\ref{fig:comp-movement}. Dividing \([0,1]\) into \(R=4\) equal subintervals
of duration \(1/4\), the \emph{spread} schedule rises by one unit during the
first and fourth subintervals and stays flat in between, with discrete height
increments
\[
 a = (a_r)_{r=0}^{3} = (1, 0, 0, 1),
\]
whereas the \emph{central} schedule places both unit rises back to back in the
middle two subintervals, with increments
\[
 b = (b_r)_{r=0}^{3} = (0, 1, 1, 0).
\]
We write \(z=(z_r)_{r=0}^{3}\) for either scalar increment sequence \(a\) or \(b\).
Both trajectories start at height \(0\), finish at height \(2\), and have time
average \(1\) over \([0,1]\). Suppose we wish to represent discrete sequences so
that a linear classifier operating on low-depth signature features can tell
such schedules apart. Seeing precisely why an ordinary depth-two signature
merges \(a\) and \(b\)---even after a time coordinate is appended---motivates the
polynomial increment lift developed below.

\begin{figure}[ht]
\centering
\includegraphics[width=.78\linewidth]{evidence/natops/movement_schedules.pdf}
\caption{Two four-step schedules on \([0,1]\) with identical initial height
\(0\), final height \(2\), and mean height \(1\) (\cite[Fig.~1]{mainpaper}).
The spread schedule \(a=(1,0,0,1)\) separates its unit rises across the
interval; the central schedule \(b=(0,1,1,0)\) places them together.}
\label{fig:comp-movement}
\end{figure}

\subsection{Piecewise-linear paths, signatures, and depth-two coordinates}
\label{app:general-signature}\label{app:meanings}
We recall the classical signature and Chen identities~\cite{chen} in the
piecewise-linear setting used here.
Let \(z_0, \ldots, z_{R-1} \in \R^d\), with \(R \geq 2\), be an ordered sequence
of vectors that will serve as path increments. Accumulating prefix sums
\[
 y_0 = 0 \in \R^d, \qquad y_k = \sum_{r=0}^{k-1} z_r \quad (1 \leq k \leq R),
\]
and interpolating linearly on each unit parameter interval \([r, r+1]\),
\[
 y(t) = y_r + (t - r) z_r \qquad (t \in [r, r+1],\ 0 \leq r < R),
\]
defines a continuous piecewise-linear path \(y:[0, R] \to \R^d\) with derivative
\(\dot y(t) = z_r\) on \((r, r+1)\).

\begin{definition}[Truncated path signature]\label{def:signature}
For an integer truncation depth \(m \geq 1\), the \emph{signature} of the
piecewise-linear path \(y:[0,R]\to\R^d\) through depth \(m\) is the collection
of iterated integrals
\begin{equation}\label{eq:app-general-signature}
 S_{\leq m}(y) = \bigl(1, S^1(y), \ldots, S^m(y)\bigr), \qquad
 S^{i_1 \cdots i_k}(y) = \int_{0 < t_1 < \cdots < t_k < R}
 dy^{i_1}_{t_1} \cdots dy^{i_k}_{t_k},
\end{equation}
where \(S^k(y) \in (\R^d)^{\otimes k}\) has entries \(S^{i_1 \cdots i_k}(y)\) for
\(1 \leq i_1, \ldots, i_k \leq d\).
\end{definition}

Because \(y\) is piecewise smooth, the Riemann--Stieltjes integrals
in~\eqref{eq:app-general-signature} reduce to ordinary Riemann integrals and
are invariant under orientation-preserving reparametrizations of \([0,R]\). For a
single linear step \(v \in \R^d\) traversed at constant velocity over a unit
interval, the ordered \(k\)-simplex \(\{0 < t_1 < \cdots < t_k < 1\}\) has
volume \(1/k!\), so the segment's truncated signature is the truncated tensor
exponential \(\exp_\otimes(v) = \sum_{k=0}^m v^{\otimes k} / k!\). By Chen's
identity, concatenating the \(R\) linear segments multiplies their truncated
tensor exponentials in chronological order.

Specializing to depth \(m=2\), each step \(z_r\) has signature
\((1, z_r, \tfrac12 z_r \otimes z_r)\), and inductively applying Chen's identity
across \(r = 0, \ldots, R-1\) yields explicit discrete formulas for the
first-level vector \(s = S^1(y) \in \R^d\) and second-level matrix
\(B = S^2(y) \in \R^{d \times d}\):
\begin{equation}\label{eq:app-depth-two}
 s = \sum_{r=0}^{R-1} z_r, \qquad
 B = \sum_{0 \leq r < t < R} z_r \otimes z_t
   + \frac{1}{2}\sum_{r=0}^{R-1} z_r \otimes z_r, \qquad
 J = \frac{1}{2}(B - B^\top)
   = \frac{1}{2}\sum_{0 \leq r < t < R} (z_r \otimes z_t - z_t \otimes z_r).
\end{equation}
The endpoint \(s\) records total net displacement. Adding~\eqref{eq:app-depth-two}
to its transpose gives the shuffle relation \(B + B^\top = s s^\top\), which
splits \(B\) into a symmetric endpoint product and an alternating part:
\[
 B_{ij} = \frac{1}{2} s_i s_j + J_{ij}.
\]
Geometrically, for each coordinate pair \(i < j\), \(J_{ij}\) is the L\'evy
signed area swept out in the \((y^i, y^j)\)-plane by the path closed with the
straight chord from \((y^i_R, y^j_R) = (s_i, s_j)\) back to the origin: a step
of size \(\alpha\) along axis \(i\) followed by a step of size \(\beta\) along
axis \(j\) encloses a triangle of signed area \(+\alpha\beta/2\). Together, the
\(d\) endpoints \(s_i\) and the \(\binom{d}{2}\) upper-triangular signed areas
\(J_{ij}\) (\(1 \leq i < j \leq d\)) form the \emph{depth-two logsignature} of
\(y\), capturing all order-sensitive information present at depth two.

We can now verify directly why ordinary depth-two coordinates fail to separate
the two schedules in Figure~\ref{fig:comp-movement}. A scalar height path
(\(d=1\)) has only an endpoint \(s_1 = 2\) and no off-diagonal area. Appending a
uniform clock \(T(t) = t\) on \([0,1]\) supplies timing through two-channel
increments \(\tilde z_r = (1/4, z_r) \in \R^2\). Evaluating~\eqref{eq:app-depth-two}
on both \(a = (1, 0, 0, 1)\) and \(b = (0, 1, 1, 0)\) gives the common endpoint
\(s = (1, 2)\) and the time--height signed area
\[
 J_{Tx} = \frac{1}{2}\sum_{0 \leq r < t < 4}
          \Bigl(\frac{1}{4} z_t - \frac{1}{4} z_r\Bigr)
        = \frac{1}{8}\sum_{r=0}^{3} (2r - 3) z_r = 0,
\]
because both schedules are palindromic (\(z_{3-r} = z_r\)). Consequently, both
clocked paths have the identical depth-two signature matrix
\(B = \begin{psmallmatrix} 1/2 & 1 \\ 1 & 2 \end{psmallmatrix}\): symmetric
rearrangements of rises around the midpoint cancel in \(J_{Tx}\).

\subsection{The Pascal increment lift and pairwise separation}
\label{app:movement}
Rather than relying on higher-depth iterated integrals against a single linear
clock, we can make step position visible before accumulating the path by
attaching polynomially weighted copies directly to each increment.

\begin{definition}[Pascal increment lift and Pascal-coded signature]
\label{def:pascal-lift}
Let \(z_0, \ldots, z_{R-1} \in \R^C\) with \(R \geq 2\), and choose an integer
copy count \(1 \leq P \leq R\). Define the normalized Pascal weights, lifted
increments, and accumulated vertices by
\begin{equation}\label{eq:app-pascal}
 b_p(r) = \frac{\binom{r}{p}}{\binom{R-1}{p}}, \qquad
 v_r = \bigl(b_p(r) z_r\bigr)_{p=0}^{P-1} \in \R^{PC}, \qquad
 y_k = \sum_{r < k} v_r \quad (0 \leq k \leq R),
\end{equation}
where \(\binom{r}{p} = 0\) for integers \(0 \leq r < p\), and \(y_0 = 0\). Writing
\(c_p\) for copy \(p \in \{0, \ldots, P-1\}\) of input channel
\(c \in \{1, \ldots, C\}\), the coordinate \(v_r^{c_p}\) equals \(b_p(r) z_r^c\).
Interpolating the vertices \(y_0, \ldots, y_R\) linearly on unit intervals defines
the lifted path \(y:[0,R]\to\R^{PC}\). Its ordinary truncated signature
\(S_{\leq m}(y)\) is the \emph{Pascal-coded signature} of \(z\) at copy count
\(P\) and depth \(m\).
\end{definition}

Two structural properties of Definition~\ref{def:pascal-lift} deserve emphasis.
First, the weights \(b_p(r)\) multiply the \emph{increments} \(z_r\), not the
accumulated positions \(x_r = \sum_{j < r} z_j\). Multiplying positions by
\(b_p(r)\) and then taking differences would inject terms proportional to
\(\bigl(b_p(r+1) - b_p(r)\bigr) x_r\) into the path increments, mixing step
changes with running totals; Definition~\ref{def:pascal-lift} instead combines
position-weighted increments~\cite{fruits} with piecewise-linear signature
accumulation~\cite{generalised}. Second, the two integers \(P\) and \(m\) govern
independent aspects of the representation: \(b_0(r) = 1\) embeds the unweighted
increment path \(x\) into the \(p=0\) coordinate block (so projecting
\(S_{\leq m}(y)\) onto the \(p=0\) channels recovers \(S_{\leq m}(x)\)),
\(b_1(r) = r/(R-1) = u_r\) adds linear position weights in \([0,1]\), and
increasing \(P\) nests higher polynomial degrees within single steps, whereas
increasing \(m\) lengthens chronologically ordered products across steps.

For scalar increments \(z=(z_r)_{r=0}^{R-1}\), take \(C=1\) and \(P=2\).
We label the lifted coordinates by copy number \(0,1\), so
\[
 v_r=(v_r^0,v_r^1)=(z_r,u_r z_r)\in\R^2.
\]
Coordinate \(0\) is the original copy and coordinate \(1\) is the
position-weighted copy. Accumulating \(v_r\) gives a two-dimensional path
\(y=(y^0,y^1)\). Its signed area in the \((y^0,y^1)\)-plane, closed by the
endpoint chord, is \(J_{01}\), as in the main paper. The indices \(0,1\)
label coordinates, while \(r,t\) label steps. Applying~\eqref{eq:app-depth-two} gives
\begin{equation}\label{eq:app-separation}
 J_{01} = \frac{1}{2}\sum_{0 \leq r < t < R} (z_r \cdot u_t z_t - u_r z_r \cdot z_t)
        = \frac{1}{2}\sum_{0 \leq r < t < R} (u_t - u_r) z_r z_t.
\end{equation}
Because \(u_t - u_r = (t - r)/(R-1) > 0\) whenever \(r < t\), every pair of
same-sign increments contributes half its product weighted by the normalized
index separation between the two steps. When increments take both signs,
positive and negative pairs can cancel, so \(J_{01}\) measures net weighted pair
separation rather than recovering individual inter-event delays. Note also that
within-step terms (\(r = t\)) vanish identically from~\eqref{eq:app-separation},
because each lifted step \((1, u_r) z_r\) moves along a straight ray in the
lifted plane and sweeps out zero area.

\begin{example}[Separating the two hand-movement schedules at depth two]
\label{ex:two-rises}
The scalar paths in Figure~\ref{fig:comp-movement} have spread increments
\(a = (1, 0, 0, 1)\) and central increments \(b = (0, 1, 1, 0)\).
Here \(R=4\), \(P=2\), and the scalar weights \(u_r=r/(R-1)\) satisfy
\((u_0,u_1,u_2,u_3)=(0,\tfrac13,\tfrac23,1)\).
For \(z=a\) or \(z=b\), the increments \(v_r=(z_r,u_r z_r)\) define the
two-dimensional coded path with coordinates \(0\) and \(1\). Both
schedules still share the coded endpoint
\[
 s(a) = \Bigl(1 + 1,\ 0\cdot 1 + 1\cdot 1\Bigr) = (2, 1), \qquad
 s(b) = \Bigl(1 + 1,\ \frac{1}{3}\cdot 1 + \frac{2}{3}\cdot 1\Bigr) = (2, 1),
\]
reflecting their common total rise and common center of rise. (If the uniform
clock \(\Delta T_r = 1/4\) is also included and lifted in copy-major order
\((T_0, x_0, T_1, x_1)\), both schedules likewise share the four-coordinate
endpoint \((1, 2, \tfrac12, 1)\).) However, each schedule has a single nonzero
pair \(r < t\) with \(z_r = z_t = 1\), so~\eqref{eq:app-separation} yields
distinct within-channel coded areas:
\[
 J_{01}(a) = \frac{1}{2}(1 - 0)\cdot 1\cdot 1 = \frac{1}{2}, \qquad
 J_{01}(b) = \frac{1}{2}\Bigl(\frac{2}{3} - \frac{1}{3}\Bigr)\cdot 1\cdot 1
           = \frac{1}{6}.
\]
Thus a linear readout on depth-two logsignature coordinates separates \(a\) from
\(b\) using the two lifted spatial channels \((y^0, y^1)\)---the same channel
count as the uncoded clocked path \((T, x)\) whose depth-two coordinates
coincided. The separation arises because \((y^0, y^1)\) pairs height increments
against each other across time, whereas \(J_{Tx}\) in \((T, x)\) pairs each
height increment only against the uniform clock.
\end{example}

\subsection{Linear readouts and channel-interaction graphs}
Understanding how coded coordinates feed a linear classifier requires
distinguishing information content from linear span. By \(B = J + \tfrac12 s s^\top\)
and \(J = \tfrac12(B - B^\top)\), the tensor pair \((s, B)\) and logsignature pair
\((s, J)\) determine one another bijectively. However, a linear functional of
\((s, B)\) has direct access to all quadratic endpoint products \(s_i s_j\),
whereas a linear functional of \((s, J)\) does not. Notice also that \(B\) alone
cannot replace \((s, B)\): because \(B\) is homogeneous of degree two in the
increments, negating an entire path (\(v_r \mapsto -v_r\)) leaves \(B\)
unchanged while reversing \(s\).

This distinction appears concretely when comparing two input channels
\(c, e \in \{1, \ldots, C\}\) under the two-copy lift \(u_r = r/(R-1)\). Write
\(c_0, c_1\) and \(e_0, e_1\) for their constant and linearly weighted copies.
The directed cross-channel separation sum
\[
 Q_{ce} = \sum_{0 \leq r < t < R} (u_t - u_r) z_r^c z_t^e
        = B_{c_0 e_1} - B_{c_1 e_0}
\]
is an exact linear combination of depth-two tensor coordinates (their within-step
terms \(\tfrac12 u_r z_r^c z_r^e\) cancel in the difference). Rewriting \(B\) in
terms of \((s, J)\) splits \(Q_{ce}\) into a linear area combination plus a
quadratic endpoint determinant:
\[
 Q_{ce} = (J_{c_0 e_1} + J_{e_0 c_1})
        + \frac{1}{2}(s_{c_0} s_{e_1} - s_{c_1} s_{e_0}).
\]
In logsignature coordinates, the two natural linear combinations of cross-copy
areas instead isolate symmetric pair separation and midpoint-weighted order:
\begin{align*}
 J_{c_0 e_1} + J_{e_0 c_1}
   &= \frac{1}{2}\sum_{0 \leq r < t < R} (u_t - u_r)(z_r^c z_t^e + z_r^e z_t^c), \\
 J_{c_0 e_1} - J_{e_0 c_1}
   &= \frac{1}{2}\sum_{0 \leq r < t < R} (u_r + u_t)(z_r^c z_t^e - z_r^e z_t^c).
\end{align*}
When \(c = e\), the endpoint determinant \(s_{c_0} s_{c_1} - s_{c_1} s_{c_0}\)
vanishes identically and \(Q_{cc} = 2 J_{c_0 c_1}\), which is why single-channel
pair separation belongs linearly to the depth-two logsignature.

Exposing multichannel interactions through \(P\) copies carries a quadratic
dimension cost. Lifting \(C\) channels to \(d = PC\) channels produces
\(\sum_{k=1}^m d^k\) non-scalar tensor coordinates through depth \(m\)
(\(d + d^2 = PC + P^2 C^2\) at depth two) and \(d + \binom{d}{2}\) depth-two
logsignature coordinates. When kinematic or sensor structure suggests retaining
cross-channel areas only across a specified graph \(E \subseteq \binom{\{1,\ldots,C\}}{2}\)
of unordered distinct channel pairs, keeping all \(PC\) endpoints, all
\(\binom{P}{2}\) same-channel cross-copy areas per input channel, and all \(P^2\)
copy-pair areas for each edge \(\{c, e\} \in E\) yields the nominal dimension
\begin{equation}\label{eq:graph-dimension}
 d_{\mathrm{graph}}(C, P, E) = PC + P^2 |E| + C\binom{P}{2}.
\end{equation}
Taking \(|E| = \binom{C}{2}\) recovers \(PC + \binom{PC}{2}\); taking \(E = \emptyset\)
retains only endpoints and single-channel separation areas, matching the sum of
\(C\) independent scalar depth-two logsignatures,
\(C\bigl(P + \binom{P}{2}\bigr)\). More generally, partitioning the channels into
disjoint groups and computing a depth-two logsignature within each group omits
cross-group areas while retaining all within-group interactions. All such
formulas count nominal coordinates; when an observation convention introduces
constant channels (such as a unit clock or unit-sum spectrum), the number of
nonconstant or linearly independent features is smaller.

\subsection{What increments represent: movement, posture, and rhythm}
\label{app:spectral-meaning}
The algebraic formulas~\eqref{eq:app-depth-two}--\eqref{eq:app-separation} are
agnostic to the physical origin of \(z_r\). Choosing how raw measurements map to
increments determines whether endpoints and areas summarize spatial gestures,
static postures, or oscillatory spectra:

\paragraph{1. Anchored displacements for spatial trajectories.}
In landmark-based action recognition~\cite{yang,generalised}, given sequential
coordinates \(x_0, \ldots, x_R \in \R^C\), taking first differences
\(z_r = x_{r+1} - x_r\) makes the path sensitive to movement while discarding any
global translation \(x_r \mapsto x_r + c\). When absolute starting pose also
distinguishes classes, prepending an initial step \(z_{\mathrm{init}} = x_0 - 0\)
from a fixed coordinate origin anchors the path: the unweighted endpoint becomes
the final pose \(x_R\), and signed areas between \(z_{\mathrm{init}}\) and
subsequent displacements record motion relative to the initial configuration.
Section~\ref{app:natops} applies this convention to 3D body-landmark gestures.

\paragraph{2. Scaled sensor readings for persistent posture levels.}
For body-worn accelerometers and gyroscopes, differencing \(a_{r+1} - a_r\)
annihilates constant gravitational orientation and merges static postures such
as sitting, standing, and lying down. Standardizing sensor readings \(a_r \in \R^C\)
using training-set statistics \emph{without} centering individual windows and
setting
\[
 z_r = \frac{a_r}{R} \qquad (0 \leq r < R)
\]
instead builds a cumulative path whose constant-copy endpoint
\(s_0 = R^{-1}\sum_{r=0}^{R-1} a_r\) is the window mean. Static posture levels
remain directly accessible, while higher Pascal copies \(b_p(r) z_r\) supply
temporal trends and ordered cross-axis interactions across the window
(Sections~\ref{app:har} and~\ref{app:additional}).

\paragraph{3. Frequency-ordered power shares for spectral shape and spread.}
Sustained rhythmic activities such as walking or cycling concentrate variance in
narrow frequency bands. To describe this concentration independently of the
phases of the Fourier coefficients, we apply the Pascal lift to power ordered
along the frequency axis~\cite{mainpaper}. Let
\(e_0, \ldots, e_{K-1} \geq 0\) be the one-sided Fourier power of a real scalar
signal at \(K \geq 2\) strictly increasing frequencies \(f_0 < \cdots < f_{K-1}\).
Whenever total power \(E = \sum_{k=0}^{K-1} e_k\) is positive, the normalized
shares \(p_k = e_k / E\) satisfy \(p_k \geq 0\) and \(\sum_{k=0}^{K-1} p_k = 1\),
defining the law of a discrete random frequency \(F\) with \(\Pr(F = f_k) = p_k\).

Here, energy refers to squared signal amplitude. Under the Parseval normalization
used in Section~\ref{app:har}, a scalar window \(x_0,\ldots,x_{n-1}\) has total
power \(E=n^{-1}\sum_r x_r^2\): the average squared reading. Squaring counts
both positive and negative readings as signal strength. For a fixed preprocessed
window, doubling every reading quadruples \(E\) but leaves the shares \(p_k\)
unchanged. The zero-frequency power \(e_0=\bar x^2\) records the constant
(mean) component.

Treating the scalar shares \(z_k = p_k\) as increments ordered by frequency bin
\(k = 0, \ldots, K-1\) and lifting with monomial frequency weights
\((1, f_k, f_k^2)\) gives coded endpoints equal to the spectral moments
\[
 s_0 = \sum_{k=0}^{K-1} p_k = 1, \qquad
 s_1 = \sum_{k=0}^{K-1} f_k p_k = \E[F], \qquad
 s_2 = \sum_{k=0}^{K-1} f_k^2 p_k = \E[F^2],
\]
namely unit mass, the spectral centroid, and the noncentral second moment.
The spectral centroid is the energy-weighted average frequency. For example,
energy shares \(3/4\) at \(2\,\mathrm{Hz}\) and \(1/4\) at \(6\,\mathrm{Hz}\)
give a centroid of \(3\,\mathrm{Hz}\). It describes where the power is centered,
which need not be a frequency present in the signal. The main paper uses the
normalized grid \(u_k=k/(K-1)\), so its centroid \(\sum_k u_kp_k\) lies in
\([0,1]\); for equally spaced frequencies starting at zero, multiplying by
\(f_{K-1}\) converts it to the original frequency units.

For an independent copy \(G\) drawn from the same distribution \(\Pr(G = f_l) = p_l\),
the pair-separation identity~\eqref{eq:app-separation} gives the constant--linear
signed area a clean probabilistic meaning:
\begin{equation}\label{eq:spectral-separation}
 J^{(1,f)}_{01} = \frac{1}{2}\sum_{0 \leq k < l < K} (f_l - f_k) p_k p_l
                = \frac{1}{4}\sum_{k=0}^{K-1}\sum_{l=0}^{K-1} |f_l - f_k| p_k p_l
                = \frac{1}{4}\E|F - G|,
\end{equation}
since each unordered pair \(\{k, l\}\) with \(k \neq l\) appears twice in the
double sum and diagonal terms vanish. Thus \(4 J^{(1,f)}_{01}\) is the classical
\emph{Gini mean difference} of the normalized power spectrum, using the full
mean-absolute-difference normalization. Its area interpretation is the
classical Gini--Lorenz relation~\cite[Sec.~2]{gini}. Shifting every frequency by a
constant leaves \(J^{(1,f)}_{01}\) unchanged, while splitting a unit spectral line
at \(f_*\) into two equal sidebands of mass \(1/2\) at \(f_* - \delta\) and
\(f_* + \delta\) keeps the centroid \(\E[F] = f_*\) fixed and increases
\(J^{(1,f)}_{01}\) from \(0\) to \(\delta/4\). (This captures spread beyond the
centroid without implying statistical independence of \(s_1\) and \(J^{(1,f)}_{01}\)
across signals.)

Because normalizing by \(E\) removes overall signal strength, a spectral shape
logsignature must be paired with separate amplitude summaries whenever vibration
intensity matters. Writing \(a_k = \sqrt{e_k}\) for the unnormalized bin
amplitudes, retaining their mean \(\bar a = K^{-1}\sum_{k=0}^{K-1} a_k\) and
population standard deviation \(\sigma_a = \sqrt{K^{-1}\sum_{k=0}^{K-1}(a_k - \bar a)^2}\)
preserves total spectral power through the exact identity
\[
 E = \sum_{k=0}^{K-1} a_k^2 = K(\bar a^2 + \sigma_a^2),
\]
which equals the time-domain mean square under the Parseval-normalized one-sided
Fourier transform specified in Section~\ref{app:har}.

\FloatBarrier\Needspace{10\baselineskip}
\section{Polynomial bases, shift equivariance, and Volterra structure}
\label{app:offsets}

Polynomial position weighting does not single out one coordinate basis on the
space of polynomials of degree below \(P\). Separating the subspace of position
summaries retained by the lift from the choice of basis clarifies three points
needed later: why \(P\) copies interpolate between unweighted signatures and full
sequence reconstruction, why positive-diagonal QR orthonormalization produces a
canonical discrete orthogonal basis for HAR and DSADS, and why the binomial
Pascal basis itself admits an exact triangular shift law and a position-only
Volterra representation.
We record standard polynomial basis and translation identities in the notation
needed for these coded signatures.

\subsection{From exact slot labels to polynomial moments}
An uncompressed representation of discrete positions would keep every time step
in its own dedicated block of \(C\) channels. Let \(e_0, \ldots, e_{R-1}\) denote
the standard basis of \(\R^R\). Assigning \(e_r \otimes z_r \in \R^R \otimes \R^C \cong \R^{RC}\)
to step \(r\) records both the increment \(z_r\) and its exact index \(r\), at
the cost of \(RC\) channels. Any linear code \(\mathcal A : \R^R \to \R^P\)
compresses these \(R\)-dimensional slot indicators into \(P\)-dimensional code
vectors \(\mathcal A e_r = b(r) \in \R^P\):
\begin{equation}\label{eq:offset-code}
 \mathcal A e_r = b(r), \qquad
 (\mathcal A \otimes I_C)(e_r \otimes z_r) = b(r) \otimes z_r = v_r,
\end{equation}
where we order coordinates copy-major so \(b(r) \otimes z_r = (b_0(r) z_r, \ldots, b_{P-1}(r) z_r)^\top\).
Position is compressed linearly within each step without shifting the sequence
in time.

Let \(S_{\mathrm{exact}}(z)\) denote the signature of the piecewise-linear path
in \(\R^{RC}\) accumulated from the exact-slot increments \(e_r \otimes z_r\).
Because the accumulated coded path satisfies \(y(t) = (\mathcal A \otimes I_C) x_{\mathrm{exact}}(t)\),
multilinearity of the iterated integrals~\eqref{eq:app-general-signature} makes
each level \(S^k_{\mathcal A}(z) = S^k(y)\) a linear image of the exact-slot
signature:
\begin{equation}\label{eq:offset-projection}
 S^k_{\mathcal A}(z) = (\mathcal A \otimes I_C)^{\otimes k} S^k_{\mathrm{exact}}(z),
 \qquad 1 \leq k \leq m.
\end{equation}

For the normalized Pascal weights \(b_p(r) = \binom{r}{p} / \binom{R-1}{p}\) of
Definition~\ref{def:pascal-lift}, the numerators \(\binom{r}{p} = r(r-1)\cdots(r-p+1)/p!\)
(\(0 \leq p < P\)) form the Newton falling-factorial basis for polynomials of
degree at most \(P-1\). Because a nonzero polynomial of degree below \(P\) can
have at most \(P-1\) roots, every \(P \times P\) minor of the \(P \times R\)
code matrix \(\mathcal A\) is nonzero. At \(P = 1\), \(b_0(r) = 1\) collapses all
slot labels and recovers the ordinary signature of \(z\). For \(1 < P < R\), the
first signature level retains \(P\) polynomial moments per input channel,
\begin{equation}\label{eq:offset-moments}
 s_{c_p} = \sum_{r=0}^{R-1} b_p(r) z_r^c,
\end{equation}
allowing a linear readout to express any degree-\((P-1)\) temporal weighting
\(\sum_{p=0}^{P-1} \beta_{pc} b_p(r)\), while level two adds polynomially
weighted ordered pairs. At the full rank \(P = R\), \(\mathcal A \in \R^{R \times R}\)
is invertible---indeed lower triangular up to column indices, with explicit
binomial inverse
\[
 z_r^c = \sum_{p=r}^{R-1} (-1)^{p-r} \binom{p}{r}\binom{R-1}{p} s_{c_p}
 \qquad (0 \leq r < R)
\]
---so the first-level coded endpoints alone reconstruct the entire sequence
\(z_0, \ldots, z_{R-1}\) linearly.

\subsection{Basis equivalence and the canonical QR orthogonal basis}\label{app:qr}
On the normalized grid \(u_r = r/(R-1)\) (\(R \geq 2\)), each Pascal weight
\(b_p(r)\) is a degree-\(p\) polynomial in \(u_r\) with strictly positive
leading coefficient \((R-1)^p / \bigl(p!\binom{R-1}{p}\bigr)\). For the first
three copies (\(R \geq 3\)),
\[
 b_0(r) = 1, \qquad
 b_1(r) = u_r, \qquad
 b_2(r) = \frac{(R-1)u_r^2 - u_r}{R-2}.
\]
Although \(b_2(r)\) differs from the pure square \(u_r^2\) by a linear term,
the functions \(b_0, \ldots, b_{P-1}\) span the exact same \(P\)-dimensional
polynomial space on \(\{0, \ldots, R-1\}\) as the ascending monomials
\(1, u_r, \ldots, u_r^{P-1}\).

More generally, if \(\widetilde b(r) = L b(r)\) for an invertible transition
matrix \(L \in \R^{P \times P}\), then in copy-major channel order the lifted
path transforms by \(\widetilde y(t) = F y(t)\) with \(F = L \otimes I_C\), and
multilinearity gives
\begin{equation}\label{eq:basis-naturality}
 S^k(\widetilde y) = F^{\otimes k} S^k(y) \quad (1 \leq k \leq m), \qquad
 \widetilde s = F s, \qquad
 \widetilde J = F J F^\top.
\end{equation}
Thus any invertible change of polynomial basis induces an invertible linear
transformation on each complete tensor signature level and on the complete
depth-two logsignature \((s, J)\). Two caveats accompany~\eqref{eq:basis-naturality}:
an invertible basis change need not map a restricted coordinate subset into
itself, and a change of basis can alter the feature Gram matrix governing ridge
regression, even after columnwise standardization
(Subsection~\ref{app:fixed-endpoints}).

Sampled Pascal or monomial columns become increasingly collinear on \([0,1]\) as
\(P\) grows. For the UCI HAR and DSADS representations we therefore orthonormalize
the sampled weights across sequence positions. Let \(A \in \R^{R \times P}\) have
unnormalized binomial entries \(A_{rp} = \binom{r}{p}\) and set
\(D = \operatorname{diag}_p\binom{R-1}{p}\). A thin QR factorization uniquely
writes
\[
 A = Q H, \qquad Q^\top Q = I_P, \qquad H_{pp} > 0 \quad (0 \leq p < P),
\]
where \(Q \in \R^{R \times P}\) has orthonormal columns and \(H \in \R^{P \times P}\)
is upper triangular. Requiring a strictly positive diagonal \(H_{pp} > 0\)
removes the sign ambiguity of the orthonormal columns (the entries
\(q_p(r) = Q_{rp}\) themselves change sign across \(r\) for \(p \geq 1\)). Since
row \(r\) of \(A\) is \(b(r)^\top D\), the orthonormal weight vector satisfies
\(q(r) = (D H^{-1})^\top b(r)\), an explicit invertible basis change of the
form~\eqref{eq:basis-naturality}.

The following is the usual uniqueness of discrete polynomial orthogonalization,
expressed through positive-diagonal QR; see~\cite{discreteqr} for the connection
between discrete orthogonal polynomials and Vandermonde QR factorizations.
\begin{proposition}[Canonical discrete orthogonal basis under positive-diagonal QR]
\label{prop:qr-uniqueness}
For any \(R \geq 2\) and \(1 \leq P \leq R\), let \(U \in \R^{R \times P}\) have
ascending monomial columns \(U_{rp} = u_r^p\) with \(u_r = r/(R-1)\). In exact
arithmetic, the positive-diagonal thin QR factorizations of \(A\), of the
normalized Pascal matrix \(W = A D^{-1}\), and of \(U\) all have the identical
orthonormal factor \(Q\).
\end{proposition}
\begin{proof}
Because \(b_p(r)\) is a degree-\(p\) polynomial in \(u_r\) with positive leading
coefficient, column \(p\) of \(A\) is a linear combination of columns
\(0, \ldots, p\) of \(U\) with a strictly positive coefficient on \(u_r^p\).
Hence \(A = U R_0\) for an upper-triangular matrix \(R_0 \in \R^{P \times P}\)
with positive diagonal entries. Substituting \(A = Q H\) gives
\(U = Q (H R_0^{-1})\) and \(W = Q (H D^{-1})\). Since upper-triangular matrices
with positive diagonals form a group under multiplication and inversion, both
\(H R_0^{-1}\) and \(H D^{-1}\) are upper triangular with positive diagonals.
Uniqueness of the thin QR factorization with positive diagonal entries implies
that all three matrices share the orthonormal factor \(Q\).
\end{proof}

Proposition~\ref{prop:qr-uniqueness} identifies the columns of \(Q\) as the
normalized discrete Legendre polynomials on \(\{0, \ldots, R-1\}\). Applying
Gram--Schmidt to the first two columns \(1\) and \(u_r\) gives their closed form:
\begin{equation}\label{eq:qr-first-two}
 q_0(r) = \frac{1}{\sqrt R}, \qquad
 q_1(r) = \frac{u_r - 1/2}{d_R}, \qquad
 d_R = \sqrt{\sum_{r=0}^{R-1} \Bigl(u_r - \frac{1}{2}\Bigr)^2}
     = \sqrt{\frac{R(R+1)}{12(R-1)}}.
\end{equation}
Under scaled sensor readings \(z_r = a_r / R\), the \(p=0\) endpoint
\(\sum_{r=0}^{R-1} q_0(r) z_r = \bar a / \sqrt R\) preserves the window mean up
to \(R^{-1/2}\). (Ordering channels by sensor axis first and QR copy second, as
in the HAR and DSADS code, simply conjugates \(F\) in~\eqref{eq:basis-naturality}
by a fixed permutation matrix.)

For a unit-sum power spectrum \(p_0, \ldots, p_{K-1}\) on \(K\) equally spaced
frequency bins \(u_k = k/(K-1)\),~\eqref{eq:qr-first-two} also determines the
exact scale of the QR spectral spread coordinate. Because \(J\) is alternating,
subtracting \(\tfrac12 q_0(k)\sqrt K = 1/2\) from \(u_k\) adds a multiple of the
constant column \(q_0\) to \(q_1\) and contributes zero signed area. Hence the
stored HAR spectral area \(\sqrt 2\,J^Q_{01} = (B^Q_{01} - B^Q_{10})/\sqrt 2\) is
an exact positive multiple of the Gini mean difference~\eqref{eq:spectral-separation}:
\begin{equation}\label{eq:spectral-qr-scale}
 \sqrt 2\,J^Q_{01}
   = \frac{\sqrt 2}{\sqrt K\,d_K} J^{(1,u)}_{01}
   = \frac{\sqrt 2}{4\sqrt K\,d_K}\,\E|U - V|,
 \qquad d_K = \sqrt{\frac{K(K+1)}{12(K-1)}},
\end{equation}
where \(U, V\) are independent draws on \(\{u_0, \ldots, u_{K-1}\}\) with
\(\Pr(U = u_k) = p_k\). For an \(N\)-sample time window, physical frequency as a
fraction of the Nyquist rate is \(F = 2k/N = \frac{2(K-1)}{N} U\); thus \(F = U\)
for even \(N = 128\) (\(K = 65\)), and \(F = \frac{126}{127} U\) for
first-difference signals of odd length \(N = 127\) (\(K = 64\)).

\subsection{Shift equivariance and concatenation of locally coded segments}
If positive-diagonal QR supplies the best-conditioned basis on a fixed grid, the
distinctive algebraic virtue of the binomial basis is how it transforms under
index translations: when separately coded adjacent segments are concatenated, the
later segment's local origin must be shifted forward, and in binomial coordinates
that shift is unit lower triangular with integer entries.

Let \(a(r) = \bigl(\binom{r}{0}, \ldots, \binom{r}{P-1}\bigr)^\top \in \R^P\),
where for arbitrary integers \(h \in \Z\) we use the falling-factorial polynomial
\(\binom{h}{p} = h(h-1)\cdots(h-p+1)/p!\) for \(p \geq 1\) and \(\binom{h}{0} = 1\).
(On nonnegative integers \(r \geq 0\), this coincides with setting
\(\binom{r}{p} = 0\) for \(r < p\) in~\eqref{eq:app-pascal} while extending
smoothly to negative shifts.) By Vandermonde's polynomial identity
\(\binom{r+h}{p} = \sum_{q=0}^p \binom{h}{p-q}\binom{r}{q}\), every integer shift
\(h \in \Z\) acts linearly on \(a(r)\):
\begin{equation}\label{eq:offset-shift}
 a(r + h) = T_h a(r), \qquad
 (T_h)_{pq} = \begin{cases}
   \binom{h}{p-q}, & 0 \leq q \leq p < P,\\
   0, & 0 \leq p < q < P.
 \end{cases}
\end{equation}
The matrices \(\{T_h\}_{h \in \Z}\) are unit lower triangular and form an exact
one-parameter abelian group: \(T_0 = I_P\), \(T_h T_k = T_{h+k}\), and
\(T_h^{-1} = T_{-h}\). For \(P = 3\), advancing the step index by \(h = 1\)
applies
\[
 T_1 = \begin{pmatrix} 1 & 0 & 0 \\ 1 & 1 & 0 \\ 0 & 1 & 1 \end{pmatrix},
 \qquad
 T_1 a(2) = \begin{pmatrix} 1 & 0 & 0 \\ 1 & 1 & 0 \\ 0 & 1 & 1 \end{pmatrix}
            \begin{pmatrix} 1 \\ 2 \\ 1 \end{pmatrix}
          = \begin{pmatrix} 1 \\ 3 \\ 3 \end{pmatrix} = a(3).
\]
Because translating the index before coding equals applying \(T_h\) after coding,
the Pascal lift is \emph{shift equivariant}. In normalized Pascal coordinates
\(b(r) = D^{-1} a(r)\) with \(D = \operatorname{diag}_p\binom{R-1}{p}\), or after
a further invertible basis map \(L\) such as QR, the shift operator conjugates to
\(D^{-1} T_h D\) or \(L D^{-1} T_h D L^{-1}\).

Combining shift equivariance with Chen's identity~\cite{chen} yields a
concatenation rule for locally coded segments without revisiting raw increments.
For any segment \(z = (z_0, \ldots, z_{\ell-1})\) of
length \(\ell \geq 1\) coded from its own local origin \(j = 0, \ldots, \ell-1\)
in unnormalized Pascal coordinates, its truncated depth-\(m\) tensor signature is
\begin{equation}\label{eq:reset-signature}
 G(z) = \prod_{j=0}^{\ell-1} \exp_\otimes\bigl(a(j) \otimes z_j\bigr).
\end{equation}
Let \(\Theta_h\) denote the truncated tensor-algebra automorphism that applies
\(T_h \otimes I_C\) to each tensor factor:
\(\Theta_h|_{(\R^{PC})^{\otimes k}} = (T_h \otimes I_C)^{\otimes k}\). Since
\((T_\ell \otimes I_C)\bigl(a(j) \otimes z_j\bigr) = a(\ell + j) \otimes z_j\)
and \(\exp_\otimes\) commutes with linear maps, concatenating \(z\) with a second
locally coded segment \(w = (w_0, \ldots, w_{q-1})\) satisfies the
\emph{twisted Chen identity}
\begin{equation}\label{eq:reset-chen}
 G(z \mathbin\Vert w) = G(z)\,\Theta_\ell\bigl(G(w)\bigr).
\end{equation}
Associativity across multiple blocks follows immediately from the group law
\(\Theta_\ell \Theta_q = \Theta_{\ell+q}\) and associativity of the truncated
tensor product. If segments of different lengths \(\ell\) were normalized
independently by their own denominators \(\binom{\ell-1}{q}\) (using
\(K = \min(P, \ell)\) local copies so even a singleton step \(\ell = 1\) avoids
division by zero), applying the \(P \times K\) rebasing map
\(H_{pq} = \binom{h}{p-q}\binom{\ell-1}{q}/\binom{R-1}{p}\) (\(q \leq p\)) to a
block starting at global index \(h\) embeds its local signature directly into the
global normalized signature \(S_{\leq m}(y)\). Under uniform sampling, integer
index shifts correspond to physical time translations.

\begin{remark}[Online sliding-window updates and comparison with SigAdapt]
\label{app:sigadapt}
Equation~\eqref{eq:reset-chen} also clarifies how representation choices differ
from online readout adaptation, such as SigAdapt~\cite{sigadapt} for sliding-window
time-series forecasting. For each look-back length, SigAdapt computes ordinary
tensor signatures (at depths four to six in its reported experiments) of
time- and basepoint-augmented history and covariate paths, fits a joint ridge
model, updates multipliers of its per-level signature contributions online via a
Kalman filter, and aggregates across window lengths. Its sliding-window signature
updates drop the oldest step and append the newest step via Chen's identity and basepoint
correction. Performing an analogous sliding update on locally indexed Pascal
signatures additionally applies \(\Theta_{-1}\) after left-multiplying by
\(\exp_\otimes\bigl(-a(0)\otimes z_0\bigr)\) so all remaining step indices shift
down by one before the new step is multiplied on the right at \(a(R-1)\).
Polynomial coding and online state-space adaptation address orthogonal parts of a
modelling pipeline and have not been combined here; moreover, algebraic update
identities for raw signature blocks do not determine end-to-end cost when a
representation also recomputes windowed Fourier transforms or lagged geometric
streams (Sections~\ref{app:har} and~\ref{app:compute}).
\end{remark}

\subsection{Nilpotent generator and Volterra signature specialization}
\label{app:nilpotent}
Because \(\{T_h\}_{h \in \Z}\) is a unipotent lower-triangular group, the Pascal
weights arise from a nilpotent matrix exponential, and the entire coded signature
embeds into the Volterra signature framework~\cite{volterra}. It is instructive
to derive both descriptions and keep their scopes distinct: the first is an
identity for single-step weights, whereas the second identifies iterated
integrals across all signature depths.

\begin{proposition}[Nilpotent generator of Pascal weights]\label{prop:nilpotent}
Let \(N \in \R^{P \times P}\) have ones on its first subdiagonal
(\(N_{p, p-1} = 1\) for \(1 \leq p < P\)) and zeros elsewhere, let \(e_0\) be the
first standard basis vector of \(\R^P\), and define the finite matrix logarithm
\[
 H = \log(I_P + N) = \sum_{j=1}^{P-1} \frac{(-1)^{j+1}}{j} N^j.
\]
Then \(T_h = e^{hH}\) and \(a(r) = e^{rH} e_0\) for every integer \(h, r \in \Z\).
In particular, relative to any fixed terminal index \(T \in \Z\),
\begin{equation}\label{eq:nilpotent-endpoint}
 e^{-(T - r) H} e_0 = e^{-T H} a(r).
\end{equation}
\end{proposition}
\begin{proof}
Because \(N^P = 0\), the power series for \(\log(I_P + N)\) and \(\exp(hH)\)
truncate at degree \(P - 1\), and formal power-series composition holds exactly:
\(e^H = I_P + N\) and \(e^{hH} = (I_P + N)^h\). By the generalized binomial
theorem for nilpotent perturbations,
\[
 (I_P + N)^h = \sum_{k=0}^{P-1} \binom{h}{k} N^k \qquad (h \in \Z).
\]
Since \((N^k)_{pq} = \mathbf 1_{\{p - q = k\}}\), entry \((p, q)\) of
\((I_P + N)^h\) equals \(\binom{h}{p-q}\) for \(q \leq p\) and \(0\) otherwise,
which is \(T_h\) in~\eqref{eq:offset-shift}. Evaluating \(T_r e_0\) selects column
\(q = 0\) and gives \(a(r)\); left-multiplying by \(e^{-TH}\) yields~\eqref{eq:nilpotent-endpoint}.
\end{proof}

Polynomial memory from nilpotent generators is treated within exponentially
weighted signatures (EWS) in~\cite[Prop.~2]{ews}. EWS also identifies its
fixed-horizon iterated integrals with the ordinary signature of a reweighted
path~\cite[Prop.~3]{ews}. For the present discrete lift, holding the clock at
index \(r\) during signal increment \(z_r\) makes the terminal-weighted increments
in~\eqref{eq:nilpotent-endpoint} a fixed linear transform of the Pascal increments;
the corresponding signatures then transform level by level. This comparison
uses that clock convention, rather than simultaneous linear clock and signal
interpolation.

To place the coded signature within the Volterra signature~\cite{volterra} at
every truncation level, recall that a Volterra signature replaces ordinary path
differentials by integrals against a measurable matrix-valued kernel \(K(t, s)\)
satisfying \(\sup_{t \in [0, R]} \int_0^t \|K(t, s)\|\,ds < \infty\). Let
\(x:[0, R]\to\R^C\) be the piecewise-linear path with increments \(z_r\) on
\([r, r+1]\), and define the piecewise-constant, position-only kernel
\[
 A(s) = b(r) \otimes I_C \quad (r \leq s < r + 1,\ 0 \leq r < R), \qquad
 K(t, s) = A(s) \quad (0 \leq s \leq t \leq R),
\]
with \(A(R) = b(R-1) \otimes I_C\). Because \(K\) is bounded on \([0, R]^2\), it
is Volterra admissible, and integrating \(A(s)\,dx_s\) on \([0, t]\) produces the
exact piecewise-linear lifted path \(y_t = \int_0^t A(s)\,dx_s\) of
Definition~\ref{def:pascal-lift}. At signature level \(k \geq 1\), the Volterra
iterated integral~\cite[Def.~2.14]{volterra} over the simplex
\(0 < r_1 < \cdots < r_k < t\) with terminal horizon \(r_{k+1} = \tau \geq t\)
forms the tensor product of factors \(K(r_{\ell+1}, r_\ell)\,dx_{r_\ell}\).
Because \(K(r_{\ell+1}, r_\ell) = A(r_\ell)\) is independent of its first time
argument \(r_{\ell+1}\), the \(\ell\)th differential simplifies to
\(A(r_\ell)\,dx_{r_\ell} = dy_{r_\ell}\), yielding exact equality at all levels:
\[
 \VSig(x; K)^{\tau}_{0, t} = \Sig(y)_{0, t} \qquad (0 \leq t \leq \tau \leq R).
\]
Subsequent rescaling or QR orthonormalization applies a fixed linear map
\(F^{\otimes k}\) to each level. Because \(K(t, s)\) depends only on \(s\), this
specialization preserves the classical Chen concatenation and shuffle identities
of ordinary signatures.

By contrast, while~\eqref{eq:nilpotent-endpoint} shows that weighting discrete
increments by the terminal-lag exponential \(e^{-(T-r)H} e_0\) differs from
\(a(r) \otimes z_r\) only by the fixed invertible transformation
\(e^{-TH} \otimes I_C\), that relation does \emph{not} identify higher levels of
the Pascal-coded signature with a convolution Volterra signature
\(K(t, s) = \widetilde K(t - s)\). In a genuine two-time convolution kernel, the
\(\ell\)th Volterra factor \(K(r_{\ell+1}, r_\ell)\) weights step \(r_\ell\) by
its lag relative to the \emph{next} integration variable \(r_{\ell+1}\), rather
than referencing all factors to a single terminal endpoint \(T\). This distinction
between fixed-terminal EWS and convolution Volterra weighting is also discussed
in~\cite[App.~D.2]{ews}.

\FloatBarrier\Needspace{10\baselineskip}
\section{Ordinary signature depth, staircase projection, and ridge geometry}
\label{app:depth-and-ridge}

Appending a clock or frequency coordinate to an uncoded path makes positional
information available in the full infinite signature. Two mathematical questions
connect that fact to the depth-two Pascal lift. First, when does a bounded-depth
ordinary signature recover the depth-two Pascal coordinates linearly? Second,
when two signature representations span the exact same affine linear function
space on a dataset, why can a linear ridge classifier still make different
predictions? We answer both questions in turn.

\subsection{Clocked paths at depth three and the interpolation correction}
In Example~\ref{ex:two-rises}, the uncoded clock-augmented paths \((T, x)\) for
the spread schedule \(a = (1, 0, 0, 1)\) and central schedule \(b = (0, 1, 1, 0)\)
had identical coordinates through depth two, while the two-copy Pascal area
\(J_{01}\), between the original and weighted copies with increments
\((z_r,u_r z_r)\), separated them as \(1/2\) versus \(1/6\). Moving one level higher
in the ordinary signature of \((T, x)\) on \([0, 1]\) exposes the space--time--space
iterated integral:
\[
 S^{xTx} = \int_{0 < t_1 < t_2 < t_3 < 1} dx_{t_1}\,dT_{t_2}\,dx_{t_3}
         = \int_0^1 \Bigl(\int_0^{t_2} dx_{t_1}\Bigr)\Bigl(\int_{t_2}^1 dx_{t_3}\Bigr)\,dt_2
         = \int_0^1 x_t (s_x - x_t)\,dt,
\]
since \(dT_t = dt\) and \(x_0 = 0\). With total rise \(s_x = 2\) and four uniform
clock steps of \(1/4\), integrating \(x_t(2 - x_t)\) across the four piecewise-linear
segments gives
\[
 S^{xTx}(a) = \frac{5}{6}, \qquad S^{xTx}(b) = \frac{1}{3},
\]
and the corresponding degree-three logsignature coefficients of the word \(xTx\)
are \(1/6\) for \(a\) and \(-1/3\) for \(b\). The ordinary clocked path therefore
first distinguishes the two schedules at depth three.

To see the exact relationship between \(S^{xTx}\) and the discrete separation
area \(J_{01}\) for an arbitrary scalar increment sequence
\(z=(z_r)_{r=0}^{R-1}\), with \(z_r\in\R\), over \(R \geq 2\) uniform clock steps
\(\Delta T_r = 1/R\), decompose \(S^{xTx}\) over the segment indices
\(r, m, t \in \{0, \ldots, R-1\}\) in which the three integration times
\(t_1 < t_2 < t_3\) fall:
\begin{itemize}
\item \emph{Distinct spatial segments (\(r < t\)):} If \(t_2\) falls in a
strictly intermediate segment \(r < m < t\) (of which there are \(t - r - 1\)),
the integral contributes \(\frac{1}{R} z_r z_t\) per segment. If \(t_2\) falls in
segment \(r\) with \(t_1 < t_2\), the ordered \(2\)-simplex on segment \(r\)
contributes \(\frac{1}{2R} z_r z_t\); symmetrically, \(t_2 < t_3\) in segment
\(t\) contributes another \(\frac{1}{2R} z_r z_t\). Summing these three cases
yields \(\bigl((t - r - 1) + \tfrac12 + \tfrac12\bigr)\frac{1}{R} z_r z_t = \frac{t - r}{R} z_r z_t\).
\item \emph{Within-segment triples (\(r = m = t\)):} When all three integration
variables fall inside the same segment \(r\), linear traversal of
\((\Delta T_r, z_r) = (1/R, z_r)\) over the ordered \(3\)-simplex of volume
\(1/3! = 1/6\) contributes \(\frac{1}{6R} z_r^2\).
\end{itemize}
Summing over all segments and comparing with
\(J_{01} = \frac{1}{2(R-1)}\sum_{r < t}(t - r) z_r z_t\) from~\eqref{eq:app-separation}
establishes the exact identity
\begin{equation}\label{eq:discrete-depth-correction}
 J_{01} = \frac{R}{2(R-1)} S^{xTx}
             - \frac{1}{12(R-1)} \sum_{r=0}^{R-1} z_r^2.
\end{equation}
For both \(a\) and \(b\), \(\sum_{r=0}^3 z_r^2 = 2\), so the within-segment
interpolation correction is \(2/(12 \cdot 3) = 1/18\), confirming
\(\frac{4}{6}\cdot\frac{5}{6} - \frac{1}{18} = \frac{1}{2}\) and
\(\frac{4}{6}\cdot\frac{1}{3} - \frac{1}{18} = \frac{1}{6}\).

\begin{example}[Depth-three collision under simultaneous linear interpolation]
\label{ex:depth-three-collision}
Because the within-step sum of squares \(\sum_{r=0}^{R-1} z_r^2\) is not a
function of the ordinary depth-three signature of \((T, x)\),
identity~\eqref{eq:discrete-depth-correction} prevents any universal recovery of
\(J_{01}\) from clocked depth-three coordinates. Explicitly, on the
four-step uniform clock \(\Delta T_r = 1/4\), consider the two scalar increment
sequences
\[
 z^{(1)} = (2, 1, 1, 1), \qquad
 z^{(2)} = \Bigl(\frac{17}{8},\, \frac{5}{8},\, \frac{11}{8},\, \frac{7}{8}\Bigr).
\]
Direct evaluation of the concatenated segment exponentials shows that the clocked
paths of \(z^{(1)}\) and \(z^{(2)}\) have \emph{identical} ordinary tensor
signatures at every word of length zero, one, two, and three (all \(15\) tensor
coordinates coincide, with \(s = (1, 5)\) and \(S^{xTx} = 103/24\)). Yet their sum
of squared increments differs (\(7\) versus \(121/16\)), so~\eqref{eq:discrete-depth-correction}
assigns them distinct depth-two Pascal areas:
\[
 J_{01}\bigl(z^{(1)}\bigr) = \frac{8}{3} = \frac{512}{192}, \qquad
 J_{01}\bigl(z^{(2)}\bigr) = \frac{509}{192}.
\]
Conversely, discrete index coding depends on how steps are segmented and
anchored: prepending an idle zero step \((0, 0)\) to \(a\) and \(b\) leaves the
geometric curve \((T, x)\) and all of its ordinary signature coordinates
unchanged, while changing the five-step Pascal grid to \(u_r = r/4\) and
rescaling their coded areas from \((1/2, 1/6)\) to \((3/8, 1/8)\).
\end{example}

\subsection{Exact scalar projection from a frequency--mass staircase}
\label{app:staircase-projection}
The discrepancy in~\eqref{eq:discrete-depth-correction} arose because \(T\) and
\(x\) changed simultaneously within each piecewise-linear segment. When a
one-dimensional sequence \(p_0, \ldots, p_{K-1}\) is instead embedded into \(\R^2\)
as a \emph{staircase} that holds the position coordinate constant during each
data jump and advances position only between jumps, within-step position changes
vanish identically, and every scalar Pascal-coded coordinate becomes an exact
linear combination of higher-depth staircase signature coordinates.

Let \(K \geq 3\) and \(u_k = k/(K-1)\) for \(0 \leq k < K\). Construct the
two-channel frequency--mass staircase \((F, M)\) starting at \((0, 0)\) as
follows: at bin \(k = 0\), traverse a vertical step \(\Delta M = p_0\) while
\(F = u_0 = 0\) remains fixed; for each \(k = 0, \ldots, K-2\), move horizontally
by \(\Delta F = u_{k+1} - u_k = 1/(K-1)\) with \(\Delta M = 0\), then traverse
the vertical mass step \(\Delta M = p_{k+1}\) at fixed \(F = u_{k+1}\). No
horizontal step follows the final jump \(p_{K-1}\) at \(u_{K-1} = 1\), so for a
unit-sum spectrum the endpoint is \((F_{\mathrm{end}}, M_{\mathrm{end}}) = (1, 1)\).

In signature words over the alphabet \(\{F, M\}\), let \(F^a\) denote \(a\)
consecutive \(F\) letters (with \(F^0\) the empty word). Write \(s_a\) and
\(B_{ab}\) for the first- and second-level coordinates of the scalar lift with
monomial increment copies \(u_k^a p_k\) and \(u_k^b p_k\), and set
\(J_{ab} = \tfrac12(B_{ab} - B_{ba})\). Here \(a,b\) label polynomial copies,
while \(k\) indexes frequency bins; in particular, \(J_{01}\) pairs the original
copy with the linearly weighted copy.

The following result specializes the standard polynomial one-form representation
of iterated integrals as signature functionals~\cite[Lem.~4.1]{polynomialoneforms},
applied to \(F^a\,dM\). We give explicit coefficients and the sufficient scalar
depth bound \(2P-1\).

\begin{theorem}[Exact scalar staircase projection]\label{prop:staircase-projection}
Along the \((F, M)\) staircase on \(K \geq 3\) bins, the monomial-coded endpoints
and signed areas for degrees \(0, 1, 2\) satisfy
\begin{align}
 s_0 &= S^M, & s_1 &= S^{FM}, & s_2 &= 2 S^{FFM}, \label{eq:staircase-endpoints}\\
 J_{01} &= \frac{1}{2} S^{MFM}, &&& \label{eq:staircase-area01}\\
 J_{02} &= S^{FMFM} + S^{MFFM}, &&& \label{eq:staircase-area02}\\
 J_{12} &= S^{FFMFM} + S^{FMFFM}. &&& \label{eq:staircase-area12}
\end{align}
Consequently, at any fixed \(K \geq 3\), all depth-two endpoints and signed areas
of the \(P = 3\) Pascal lift---as well as its positive-diagonal QR version
(Section~\ref{app:qr}) and \(\sqrt 2\)-scaled spectral areas---are fixed,
data-independent linear projections of the ordinary staircase tensor signature
through depth five. More generally, for \(2 \leq P \leq K\) scalar polynomial
copies of degrees below \(P\), ordinary staircase depth \(2P - 1\) suffices.
\end{theorem}
\begin{proof}
During the \(k\)th vertical segment, \(dM\) integrates to \(p_k\) at constant
frequency \(F_t = u_k\), whereas \(dM = 0\) on every horizontal segment. Since
\(F_0 = 0\), the ordered \(a\)-fold integral of \(dF\) from \(0\) up to time
\(t\) equals \(F_t^a / a!\). Integrating against \(dM_t\) therefore yields
\[
 a!\,S^{F^a M} = \int_0^T F_t^a\,dM_t = \sum_{k=0}^{K-1} u_k^a p_k = s_a,
\]
which gives~\eqref{eq:staircase-endpoints} for \(a = 0, 1, 2\).

For the second-level coordinate \(B_{ab} = \int_{0 < s < t < T} F_s^a F_t^b\,dM_s\,dM_t\)
(notice that because \(F\) is constant within each vertical mass step, integrating
\(F_s^a F_t^b\,dM_s\,dM_t\) over \(s < t\) inside the \(k\)th vertical segment
produces \(\tfrac12 u_k^{a+b} p_k^2\), matching the piecewise-linear diagonal
term in~\eqref{eq:app-depth-two}), expand \(F_t^b = \bigl(F_s + (F_t - F_s)\bigr)^b\)
by the binomial theorem:
\[
 F_s^a F_t^b = \sum_{j=0}^b \binom{b}{j} F_s^{a+b-j} (F_t - F_s)^j.
\]
For \(0 < s < t\), \(F_s^{a+b-j} / (a+b-j)!\) is the \((a+b-j)\)-fold iterated
integral of \(dF\) over \((0, s)\), and \((F_t - F_s)^j / j!\) is the \(j\)-fold
iterated integral of \(dF\) over \((s, t)\). Interleaving these \(dF\) factors
with \(dM_s\) and \(dM_t\) yields, for all integers \(a, b \geq 0\),
\begin{equation}\label{eq:staircase-general}
 B_{ab} = \sum_{j=0}^b \binom{b}{j} (a + b - j)!\,j!\,S^{F^{a+b-j} M F^j M}.
\end{equation}
Notice that the \(j = 0\) term \((a+b)!\,S^{F^{a+b} M M}\) is symmetric in \(a\)
and \(b\) and cancels upon forming \(J_{ab} = \tfrac12(B_{ab} - B_{ba})\).
Evaluating~\eqref{eq:staircase-general} at \((a, b) = (0, 1), (0, 2), (1, 2)\)
and subtracting \(B_{ba}\) gives:
\begin{align*}
 B_{01} - B_{10} &= (S^{FMM} + S^{MFM}) - S^{FMM} = S^{MFM}, \\
 B_{02} - B_{20} &= (2S^{FFMM} + 2S^{FMFM} + 2S^{MFFM}) - 2S^{FFMM}
                  = 2(S^{FMFM} + S^{MFFM}), \\
 B_{12} - B_{21} &= (6S^{FFFMM} + 4S^{FFMFM} + 2S^{FMFFM})
                  - (6S^{FFFMM} + 2S^{FFMFM}) \\
                 &= 2(S^{FFMFM} + S^{FMFFM}),
\end{align*}
proving~\eqref{eq:staircase-area01}--\eqref{eq:staircase-area12}.

Finally, any degree-\((P-1)\) polynomial basis \(c_p(u) = \sum_{a=0}^{P-1} L_{pa} u^a\)
transforms endpoints and areas linearly via \(s^c = L s\) and \(J^c = L J L^\top\)
by~\eqref{eq:basis-naturality}, including the normalized Pascal weights, the
orthonormal QR weights \(q(k) = (D H^{-1})^\top b(k)\) of Section~\ref{app:qr},
and multiplication of areas by \(\sqrt 2\). For \(0 \leq a < b \leq P - 1\), the
words in~\eqref{eq:staircase-general} have length \((a + b - j) + 1 + j + 1 = a + b + 2 \leq (P - 2) + (P - 1) + 2 = 2P - 1\),
while endpoint words \(F^a M\) have length \(a + 1 \leq P < 2P - 1\) for
\(P \geq 2\).
\end{proof}

\begin{remark}[Scope of the staircase construction]
Every word appearing in~\eqref{eq:staircase-endpoints}--\eqref{eq:staircase-general}
ends in the letter \(M\). Appending an extra horizontal step after the last mass
jump would change \(S^F\) and words ending in \(F\), while leaving every formula
in Theorem~\ref{prop:staircase-projection} intact. By contrast, replacing the
axis-aligned staircase with diagonal segments that increment \(F\) and \(M\)
simultaneously reintroduces within-bin interpolation terms analogous
to~\eqref{eq:discrete-depth-correction}. Likewise, Theorem~\ref{prop:staircase-projection}
exploits the scalar structure of \(M\) (where same-degree areas \(J_{aa}\) vanish)
and does not assert an identical depth-\((2P-1)\) projection for general
multichannel temporal paths.
\end{remark}

\subsection{Fixed endpoints, affine spans, and ridge regularization}
\label{app:fixed-endpoints}
Theorem~\ref{prop:staircase-projection} shows that the five nonconstant \(P=3\)
spectral Pascal coordinates lie inside the linear span of the staircase tensor
signature through depth five. When comparing candidate representations empirically
with a regularized linear classifier, however, matching or enlarging the linear
span is not equivalent to preserving ridge predictions.

The standard tensor exponential and logarithm identities~\cite{chen} show that,
even at depth three, fixed endpoints create exact affine equivalences between
tensor and logsignature coordinates. Writing \(\ell_1 = s\), \(\ell_2 = J\), and
\(\ell_3\) for the homogeneous components of the tensor logarithm
\(\log_\otimes S_{\leq 3}(y)\) through degree three, expanding the tensor
exponential \(S_{\leq 3}(y) = \exp_\otimes(\ell_1 + \ell_2 + \ell_3)\) gives
\begin{equation}\label{eq:fixed-endpoint-span}
 S^2 = J + \frac{1}{2} s^{\otimes 2}, \qquad
 S^3 = \ell_3 + \frac{1}{2}(s \otimes J + J \otimes s) + \frac{1}{6} s^{\otimes 3}.
\end{equation}
On a family of paths whose endpoint \(s\) is identical for every observation---such
as the unit-sum normalized \((F, M)\) staircases of Subsection~\ref{app:staircase-projection},
all of which share \(s = (1, 1)\)---the relations~\eqref{eq:fixed-endpoint-span}
and their Lie-series inverses are \emph{affine linear} in the remaining
coordinates \((J, \ell_3)\) and \((S^2, S^3)\). Consequently, through depth three,
staircase tensor features and staircase logsignature features span the exact same
affine linear function space on unit-sum spectra. (At depth four, the exponential
expansion of \(S^4\) introduces the quadratic area product \(\tfrac12 J \otimes J\),
so fixed endpoints alone no longer guarantee affine equivalence with the
depth-four logsignature.)

Why, then, can two feature maps with the same affine span produce different
classifications under a linear ridge model? Suppose \(Z \in \R^{N \times d_1}\)
and \(\widetilde Z \in \R^{N \times d_2}\) are two column-standardized training
feature matrices having the exact same column space \(\operatorname{col}(Z) = \operatorname{col}(\widetilde Z) \subseteq \{u \in \R^N : \mathbf 1^\top u = 0\}\).
For centered targets \(\widetilde Y = Y - \mathbf 1 \bar y^\top\) and ridge
penalty \(\alpha > 0\), solving
\[
 \min_{W} \|\widetilde Y - Z W\|_F^2 + \alpha \|W\|_F^2
\]
yields the standard dual ridge predictions~\cite{ridge}
\(\widehat Y = Z(Z^\top Z + \alpha I_{d_1})^{-1} Z^\top \widetilde Y = K(K + \alpha I_N)^{-1} \widetilde Y\),
determined entirely by the Gram matrix \(K = Z Z^\top\). At fixed \(\alpha\),
the two training prediction operators agree for all centered targets exactly
when \(K = \widetilde K = \widetilde Z \widetilde Z^\top\). For test features
\(Z_*\) standardized using the training statistics, the scores are
\(Z_* Z^\top (K + \alpha I_N)^{-1}\widetilde Y + \mathbf 1\bar y^\top\).
Columnwise standardization gives each nonconstant training column norm
\(\sqrt N\), while the Gram matrices can still differ.
\begin{enumerate}
\item \emph{Redundant or linearly dependent coordinates:} If \(\widetilde Z\)
repeats a standardized column \(k\) times (for example, when \(S^{FM}\) and
\(S^{MF} = 1 - S^{FM}\) both standardize to \(\pm z_1\) on unit-endpoint
staircases), that direction receives \(k\) times larger weight in
\(\widetilde Z \widetilde Z^\top\), effectively dividing its ridge penalty by
\(k\) relative to unrepeated directions.
\item \emph{Non-orthogonal basis transforms and unprojected ambient dimensions:}
Replacing \(P=3\) QR-Pascal coordinates by unorthogonalized monomials can alter
the singular values of \(Z\) without changing its column space. Fitting ridge
directly to all \(1{,}040\) depth-five staircase tensor columns can also add
directions beyond the \(100\) projected Pascal features under the same isotropic
penalty \(\alpha \|W\|_F^2\).
\end{enumerate}
Part~II's empirical comparisons---especially the staircase depth table in
Section~\ref{app:har} and the width-varying interaction controls in
Section~\ref{app:natops}---directly reflect this interplay between expressive span
and ridge metric geometry.


\clearpage
\part{Experimental Studies}
\section{Smartphone activity recognition on held-out subjects (UCI HAR)}
\label{app:har}

\subsection{Split, representation summary, and fitting protocol}
The UCI Human Activity Recognition (HAR) benchmark~\cite{har} records six
smartphone activities (walking, walking upstairs, walking downstairs, sitting,
standing, and laying) at \(50\,\mathrm{Hz}\) in fixed windows of \(R = 128\)
samples (\(2.56\,\mathrm{s}\)). We use the official subject-disjoint partition:
\(7{,}352\) training windows from \(21\) participants and \(2{,}947\) test
windows from \(9\) held-out participants. Each window supplies three-axis body
acceleration \(b_r \in \R^3\), angular velocity \(\omega_r \in \R^3\), and total
acceleration \(a_r \in \R^3\), from which we form the gravitational component
\(g_r = a_r - b_r\). No additional filtering or re-segmentation is applied.

Recognizing both static postures and dynamic gaits with a single linear readout
calls for complementary views of each window: gravitational orientation and
ordered cross-axis motion in the time domain, oscillation cadence and harmonic
spread in the frequency domain, and frame-invariant magnitude, alignment, and
short-lag persistence. The \emph{compact368} representation concatenates three
corresponding branches whose feature counts are summarized up front in
Table~\ref{tab:compact-blocks}.

\begin{table}[ht]
\centering
\begin{tabular}{@{}lrr@{}}
\toprule
\emph{compact368} component & Nominal & Retained \\
\midrule
Raw sensor \(P=3\) QR-Pascal depth-two logsignature (\(18\) channels) & 171 & 171 \\
Spectral shape (\(20\) scalar \(P=3\) QR-Pascal depth-two logsignatures) & 120 & 100 \\
Spectral scale (mean and std of unnormalized bin amplitudes, \(20\) signals) & 40 & 40 \\
Quadratic geometric streams (mean, std, skewness of \(12\) streams) & 36 & 36 \\
Instantaneous cosine alignment (mean and std of \(3\) cross-triad streams) & 6 & 6 \\
Window-level linear trends (\(q_1\)-weighted trend across all \(15\) streams) & 15 & 15 \\
\midrule
Total & 388 & 368 \\
\bottomrule
\end{tabular}
\caption{Feature breakdown of \emph{compact368} before and after applying the
training-only variance mask, which removes the \(20\) constant unit-mass
spectral endpoints.}
\label{tab:compact-blocks}
\end{table}

Every experiment on the official UCI HAR split in this section uses an identical,
strictly training-fitted preprocessing and classification pipeline:
\begin{enumerate}
\item \emph{Sensor calibration:} All per-axis training means and standard
deviations (raw and spectral branches) and per-triad root-mean-square scales
(geometric branch) are estimated across the \(7{,}352\) training windows only
and applied unchanged to test windows.
\item \emph{Constant-column variance mask:} Before fitting the classifier,
feature column \(j\) is discarded if its training population standard deviation
satisfies \(\sigma^{\mathrm{train}}_j \leq 10^{-6}\max_i |F^{\mathrm{train}}_{ij}|\).
On HAR this removes precisely the \(20\) unit-sum spectral endpoints, leaving
\(368\) columns.
\item \emph{Feature standardization and ridge readout:} Each retained column is
standardized using its training mean and training population standard deviation.
Writing \(F \in \R^{N_{\mathrm{train}} \times d}\) for the standardized feature
matrix and \(Y \in \{\pm 1\}^{N_{\mathrm{train}} \times 6}\) for one-versus-rest
target indicators (\(+1\) for the target activity, \(-1\) otherwise), we solve
the closed-form ridge problem with an unpenalized intercept \(b \in \R^6\) and
fixed penalty \(\alpha = 10\):
\begin{equation}\label{eq:har-ridge}
 \min_{b \in \R^6,\, W \in \R^{d \times 6}}
 \|Y - \mathbf 1 b^\top - F W\|_F^2 + \alpha \|W\|_F^2.
\end{equation}
Test windows are standardized with the frozen training parameters and assigned
to \(\arg\max_{c} \bigl(b_c + (F_{\mathrm{test}} W)_{ic}\bigr)\).
\end{enumerate}

\subsection{Specification of the three feature branches}

\paragraph{1. Raw sensor branch: posture levels and temporal order (171 features).}
To preserve static gravitational orientation across windows, we take the six raw
channels \((a_r, \omega_r) \in \R^6\) and standardize each axis using the global
mean and population standard deviation computed across all \(7{,}352 \times 128\)
training samples, \emph{without} centering individual windows. Denoting the
standardized reading at sample \(r \in \{0, \ldots, 127\}\) by \(x_r \in \R^6\),
we treat scaled readings \(x_r / 128\) as increments and lift them with the
\(P = 3\) orthonormal QR-Pascal weights \(q_p(r)\) of Section~\ref{app:qr}:
\[
 v_r^{(c, p)} = \frac{x_r^c\,q_p(r)}{128},
 \qquad c \in \{1, \ldots, 6\},\quad p \in \{0, 1, 2\}.
\]
Lifted coordinates are ordered by sensor axis \(c\) first and QR copy \(p\)
second, and accumulated from \(y_0 = 0\) without an added clock channel. Because
\(q_0(r) = 1/\sqrt{128}\), the constant-copy endpoint equals
\(\sum_{r=0}^{127} x_r^c / 128^{3/2} = \bar x^c / \sqrt{128}\), retaining the
window mean alongside linear and quadratic temporal moments and ordered
cross-channel areas. The depth-two logsignature of this \(18\)-channel path
contributes \(18\) endpoints and \(\binom{18}{2} = 153\) signed areas
\(J_{ij} = \tfrac12(B_{ij} - B_{ji})\), totaling \(171\) features.

\paragraph{2. Spectral branch: frequency shape and amplitude scale (140 retained features).}
To summarize rhythmic cadence and harmonic spread, we start from five
three-axis triads: \(b, g, \omega\) (length \(n = 128\)) and the
unscaled first differences \(\Delta b_r = b_{r+1} - b_r\) and
\(\Delta \omega_r = \omega_{r+1} - \omega_r\) (length \(n = 127\)). Each triad is
standardized per axis using training-set statistics as in the raw branch, after
which we append its instantaneous Euclidean norm \(\|x_r\|\) \emph{in the time
domain} to obtain \(20\) scalar signals (\(12\) of length \(128\) and \(8\) of
length \(127\)). Forming \(\|x_r\|\) prior to the Fourier transform captures
nonlinear envelope modulations; its power spectrum is distinct from the sum of
the three individual axis spectra. The power spectrum of each resulting scalar
signal is invariant under circular shifts of that signal. A circular shift of
the raw window can change its first-difference spectra, since differencing
omits the wraparound step.

For each scalar signal \(x_0, \ldots, x_{n-1}\), we compute the Parseval-scaled
one-sided discrete Fourier power spectrum over \(K = \lfloor n/2 \rfloor + 1\)
bins (\(K = 65\) for \(n = 128\), \(K = 64\) for \(n = 127\)):
\[
 \widehat x_k = \sum_{r=0}^{n-1} x_r e^{-2\pi i k r / n}, \qquad
 e_k = w_k \left|\frac{\widehat x_k}{n}\right|^2, \qquad
 E = \sum_{k=0}^{K-1} e_k = \frac{1}{n}\sum_{r=0}^{n-1} x_r^2,
\]
where \(w_k = 1\) at DC (\(k = 0\)) and, when \(n\) is even, at the Nyquist bin
(\(k = n/2\)), and \(w_k = 2\) at all other bins, including the highest bin for
odd \(n\). Normalized power shares are \(p_k = e_k / E\), and unnormalized bin
amplitudes are
\(a_k = \sqrt{e_k}\). Treating each signal's shares \(p_k\) independently as
scalar increments ordered by bin \(k = 0, \ldots, K-1\) and lifting with the
three QR weights \((q_0(k), q_1(k), q_2(k))\) gives \(3 + \binom{3}{2} = 6\)
depth-two logsignature coordinates per spectrum, with the three areas stored at
scale \(\sqrt 2\,J_{ij} = (B_{ij} - B_{ji})/\sqrt 2\) as analyzed
in~\eqref{eq:spectral-qr-scale}. Because \(\sum_{k=0}^{K-1} p_k = 1\), the first
endpoint \(q_0(0)\sum_k p_k = 1/\sqrt K\) is constant across windows and removed
by the variance mask, leaving \(20 \times 5 = 100\) nonconstant spectral-shape
features (with no cross-spectrum areas). Appending the mean and population
standard deviation of \(a_0, \ldots, a_{K-1}\) for each signal adds
\(20 \times 2 = 40\) amplitude-scale features.

\paragraph{3. Geometric branch: magnitude, persistence, and alignment (57 features).}
To capture sensor intensity, inter-triad alignment, and short-lag autocorrelation
independently of coordinate-axis sign conventions, we calibrate the uncentered
physical triads \(b, g, \omega\) by dividing each triad by its scalar root mean
square over all three axes, time samples, and training windows. Using one scalar
RMS per triad rather than per-axis centering ensures that the calibrated vectors
\(u_1(r), u_2(r), u_3(r) \in \R^3\) transform covariantly under any constant
orthogonal rotation \(O \in \mathrm{O}(3)\) applied simultaneously to all three
triads. From \(u_1, u_2, u_3\) we form \(12\) rotation-invariant quadratic
streams,
\[
 \|u_i(r)\|^2 \quad (1 \leq i \leq 3), \qquad
 \langle u_i(r), u_j(r)\rangle \quad (1 \leq i < j \leq 3), \qquad
 \langle u_i(r), u_i(r - \tau)\rangle \quad (1 \leq i \leq 3;\ \tau \in \{2, 5\}),
\]
where lags \(\tau = 2\) and \(\tau = 5\) correspond to \(40\,\mathrm{ms}\) and
\(100\,\mathrm{ms}\) (giving valid lengths \(L = 128, 126, 123\)), together with
three instantaneous cosine-similarity streams
\(\langle u_i(r), u_j(r)\rangle / (\|u_i(r)\|\|u_j(r)\|)\) (\(1 \leq i < j \leq 3\)).
Within each window, every stream \(z_0, \ldots, z_{L-1}\) contributes its mean
\(\bar z\) and population standard deviation \(\sigma_z\); each of the \(12\)
quadratic streams additionally contributes its skewness
\(L^{-1}\sum_{r=0}^{L-1} y_r^3\), where \(y_r = (z_r - \bar z)/\sigma_z\); and all
\(15\) streams contribute their linear QR-Pascal trend
\[
 T(z) = \frac{1}{L}\sum_{r=0}^{L-1} q_1(r) y_r,
\]
which is positive when standardized stream values increase toward the end of the
window. Taken alone, these \(36 + 6 + 15 = 57\) features constitute the
\emph{geometry57} representation.

\begin{remark}[Numerical safeguards and archived floating-point conventions]
To guard against degenerate denominators without altering well-conditioned
windows, per-axis training standard deviations and per-triad training RMS scales
are floored at \(10^{-8}\); spectral total power uses \(\max(E, 10^{-20})\),
with a unit-mass check after normalization that rejects zero-power inputs
(all evaluated HAR spectra summed to one to machine precision);
cosine-similarity denominators use \(\max(\|u_i\|\|u_j\|, 10^{-12})\);
and within-window standardized streams set \(y_r = 0\) identically whenever
\(\sigma_z^2 \leq 10^{-12}\) (otherwise flooring \(\sigma_z\) at \(10^{-8}\)). To
follow the archived rounding convention, raw and spectral signature tensors
are accumulated in binary64 and cast to binary32 immediately before extracting
upper-triangular areas, whereas geometric summaries and the ridge solve remain in
binary64 (\nolinkurl{reproduction/uci_har.py}).
\end{remark}

\subsection{Standalone blocks, supplied baseline, and single-block deletions}
\label{app:ablations}
On the \(2{,}947\) held-out subject test windows, \emph{compact368} achieves
\(98.03\%\) accuracy (\(58\) errors). Two complementary ablations clarify how its
components contribute to the \(\alpha = 10\) linear readout: fitting the readout
to individual blocks in isolation (Table~\ref{tab:har-standalone}), and removing
one block at a time from the full representation (Table~\ref{tab:har-deletions}).

\begin{table}[ht]
\centering
\begin{tabular}{@{}lrrr@{}}
\toprule
Standalone representation or baseline & Features & Accuracy (\%) & Errors \\
\midrule
Raw QR-Pascal logsignature only & 171 & 81.98 & 531 \\
Spectral QR-Pascal shape only (no amplitude scale) & 100 & 82.90 & 504 \\
Geometric branch (\emph{geometry57}) only & 57 & 85.14 & 438 \\
Supplied UCI HAR hand-engineered baseline & 561 & 95.72 & 126 \\
Full \emph{compact368} & 368 & 98.03 & 58 \\
\bottomrule
\end{tabular}
\caption{Standalone branch evaluations and the 561-feature benchmark baseline on
the official subject-disjoint UCI HAR split (\(2{,}947\) test windows) under the
identical \(\alpha = 10\) ridge readout.}
\label{tab:har-standalone}
\end{table}

\begin{table}[ht]
\centering
\begin{tabular}{@{}lrrrr@{}}
\toprule
Leave-one-out deletion from \emph{compact368} & Retained features & Accuracy (\%) & Errors & \(\Delta\) Errors \\
\midrule
All blocks retained (\emph{compact368}) & 368 & 98.03 & 58 & \(0\) \\
Remove raw QR-Pascal branch (\(171\)) & 197 & 96.17 & 113 & \(+55\) \\
Remove spectral QR-Pascal shape (\(100\)) & 268 & 96.78 & 95 & \(+37\) \\
Remove spectral amplitude mean/std (\(40\)) & 328 & 97.15 & 84 & \(+26\) \\
Remove cross-triad cosine mean/std (\(6\)) & 362 & 96.34 & 108 & \(+50\) \\
Remove quadratic stream mean/std/skew (\(36\)) & 332 & 97.42 & 76 & \(+18\) \\
Remove geometric/cosine linear trends (\(15\)) & 353 & 97.86 & 63 & \(+5\) \\
\bottomrule
\end{tabular}
\caption{Single-block deletions from \emph{compact368} on the official
subject-disjoint UCI HAR split (\(2{,}947\) test windows).}
\label{tab:har-deletions}
\end{table}

Taken alone, each individual branch makes between \(438\) and \(531\) errors
(\(81.98\%\)--\(85.14\%\)), well below both the 561-feature engineered baseline
(\(95.72\%\), \(126\) errors) and their concatenation (\(98.03\%\), \(58\)
errors). Conversely, Table~\ref{tab:har-deletions} shows that deleting any major
component substantially increases test errors even when all remaining blocks are
present: dropping the raw Pascal branch adds \(55\) errors, dropping the six
cross-triad cosine features adds \(50\) errors, dropping the spectral Pascal
shape coordinates adds \(37\) errors, and dropping the \(40\) unnormalized
amplitude summaries adds \(26\) errors (confirming empirically why unit-normalized
spectra require separate energy summaries). Because ridge weights are refitted
for each subset and correlated signal properties are shared across branches,
these error increments measure conditional usefulness to the linear readout
rather than isolated causal effects. Aggregate error counts are reported because
per-window subject-indexed predictions were not saved in the original
\emph{compact368} archive.

\subsection{Conventional signature replacements and staircase-projection verification}
To test what Pascal increment coding contributes beyond conventional path
signatures within the same pipeline, an archived \(187\)-configuration diagnostic
study replaced either the raw branch or the spectral branch with uncoded
clock-augmented signatures while holding the remaining \emph{compact368} features
and the \(\alpha = 10\) readout fixed.

\paragraph{Replacing the raw sensor branch.}
Keeping the \(140\) spectral and \(57\) geometric features fixed (\(197\)
columns), we replaced the \(171\)-feature raw QR-Pascal block with ordinary
tensor and logsignatures (up to \(1{,}600\) retained raw coordinates) built from
the same six standardized sensor axes. The strongest conventional replacement in
this setting used a clock-augmented, origin-anchored \emph{lead--lag} path: for
each consecutive pair of standardized readings \(x_r, x_{r+1} \in \R^6\), the
\(12\) spatial coordinates visit \((x_r, x_r)\), \((x_{r+1}, x_r)\), and
\((x_{r+1}, x_{r+1})\) in order, augmented by a uniform clock from \(0\) to \(1\)
and an initial step from the origin \(0 \in \R^{13}\). Removing the deterministic
clock endpoint and clock--clock square leaves \(13 + 13^2 - 2 = 180\) depth-two
tensor coordinates. Table~\ref{tab:ordinary-raw} compares this best observed
lead--lag control against \emph{compact368}, alongside a parallel comparison
appended to the \(561\) supplied HAR features (where the conventional control
visits each six-axis reading once with a clock and origin anchor, giving
\(7 + 7^2 - 2 = 54\) coordinates, against all \(18 + 18^2 = 342\) depth-two
tensor coordinates of the \(P = 3\) QR-Pascal lift).

\begin{table}[ht]
\centering
\begin{tabular}{@{}lrrrr@{}}
\toprule
Raw sensor branch representation & Raw & Total & Accuracy (\%) & Errors \\
\midrule
\multicolumn{5}{@{}l}{\emph{Within \emph{compact368} (fixed 140 spectral + 57 geometric features)}} \\
QR-Pascal \(P=3\), depth-two logsignature (\emph{compact368}) & 171 & 368 & 98.03 & 58 \\
Clocked anchored lead--lag tensor signature, depth 2 & 180 & 377 & 97.46 & 75 \\
\midrule
\multicolumn{5}{@{}l}{\emph{Appended to the 561 supplied UCI HAR features}} \\
QR-Pascal \(P=3\), depth-two tensor signature & 342 & 903 & 97.12 & 85 \\
Clocked anchored ordinary tensor signature, depth 2 & 54 & 615 & 97.05 & 87 \\
\bottomrule
\end{tabular}
\caption{Raw-branch replacements on the official subject-disjoint UCI HAR split
under fixed \(\alpha = 10\) ridge classification (\cite[Table~1]{mainpaper}).}
\label{tab:ordinary-raw}
\end{table}

\paragraph{Replacing the spectral branch and checking Theorem~\ref{prop:staircase-projection}.}
Holding the \(171\) raw features, \(57\) geometric features, and \(40\) spectral
amplitude summaries fixed (\(268\) columns), we replaced the \(100\)
QR-Pascal spectral-shape features with ordinary signatures of the two-channel
frequency--mass staircase \((F, M)\) defined in Subsection~\ref{app:staircase-projection},
evaluated independently on each of the \(20\) normalized spectra across depths
\(2\) through \(5\). Because every spectrum has fixed staircase endpoint
\((1, 1)\), words consisting purely of \(F\) letters or purely of \(M\) letters
are constant and discarded (removing both endpoints from the logsignatures).
Table~\ref{tab:ordinary-spectral} reports the resulting dimensions and test
scores, concluding with the explicit linear projection of the depth-five
staircase tensor onto the \(P=3\) QR-Pascal coordinates.

\begin{table}[ht]
\centering
\begin{tabular}{@{}lrrrr@{}}
\toprule
Spectral shape branch & Spectral & Total & Accuracy (\%) & Errors \\
\midrule
QR-Pascal \(P=3\), depth-two logsignature (\emph{compact368}) & 100 & 368 & 98.03 & 58 \\
Ordinary \((F,M)\) staircase logsignature, depth 2 & 20 & 288 & 97.49 & 74 \\
Ordinary \((F,M)\) staircase logsignature, depth 3 & 60 & 328 & 97.79 & 65 \\
Ordinary \((F,M)\) staircase tensor signature, depth 3 & 160 & 428 & 97.83 & 64 \\
Ordinary \((F,M)\) staircase logsignature, depth 4 & 120 & 388 & 97.83 & 64 \\
Ordinary \((F,M)\) staircase tensor signature, depth 4 & 440 & 708 & 98.00 & 59 \\
Ordinary \((F,M)\) staircase logsignature, depth 5 & 240 & 508 & 97.93 & 61 \\
Ordinary \((F,M)\) staircase tensor signature, depth 5 & 1,040 & 1,308 & 97.83 & 64 \\
\midrule
Depth-five staircase projected to \(P=3\) QR-Pascal & 100 & 368 & 98.03 & 58 \\
\bottomrule
\end{tabular}
\caption{Spectral-shape replacements within \emph{compact368} (with the other
\(268\) raw, amplitude-scale, and geometric features held fixed) on the official
UCI HAR split. Counts are retained columns, including linear dependencies among
tensor words. (The depth-two staircase tensor with \(40\) spectral columns also
achieves \(97.49\%\), \(74\) errors.)}
\label{tab:ordinary-spectral}
\end{table}

Table~\ref{tab:ordinary-spectral} illustrates three distinct theoretical points
from Section~\ref{app:depth-and-ridge}:
\begin{enumerate}
\item \emph{Expressive depth versus compactness:} Raising staircase depth from
two (\(74\) errors) to four in tensor coordinates (\(59\) errors, \(98.00\%\))
essentially matches \emph{compact368} (\(58\) errors, \(98.03\%\)---a single
test window, which is descriptively tied), using \(440\) spectral columns
(\(708\) total) in place of \(100\) (\(368\) total). Continuing to all \(1{,}040\)
unprojected depth-five tensor columns slightly increases errors to \(64\) under
the fixed isotropic penalty \(\alpha = 10\).
\item \emph{Fixed-endpoint ridge sensitivity:} At depth three, the staircase
logsignature (\(60\) spectral columns, \(65\) errors) and staircase tensor
signature (\(160\) spectral columns, \(64\) errors) span the \emph{exact same}
affine linear function space by~\eqref{eq:fixed-endpoint-span}, yet yield
slightly different ridge fits because redundant tensor words alter the singular
values of the standardized feature matrix (Subsection~\ref{app:fixed-endpoints}).
\item \emph{Numerical verification of the staircase projection.} The archived
self-check compares the projected and directly computed QR-Pascal coordinates
on synthetic spectra with absolute tolerance \(10^{-7}\) and relative tolerance
\(10^{-6}\). For the final row of Table~\ref{tab:ordinary-spectral}, the archived
run reports \emph{identical class predictions on all \(2{,}947\) test windows}
after replacing the \(100\) spectral-shape features by the projection. The
mathematical identity holds for all spectra by Theorem~\ref{prop:staircase-projection}.
\end{enumerate}
\FloatBarrier

\FloatBarrier\Needspace{10\baselineskip}
\section{Landmark trajectories and interaction graphs (NATOPS)}
\label{app:natops}

\subsection{Anchored displacement path and interaction graphs}
Whereas smartphone accelerometry relies on persistent gravitational levels,
skeletal landmark sequences describe gestures through coordinated spatial motion
from an initial posture. The NATOPS benchmark~\cite{natopsoriginal,natops} provides a clean
setting to probe which cross-channel interactions benefit from Pascal coding in a
small-sample, high-width regime: six aircraft-handling command gestures are
recorded over \(51\) frames across eight 3D body landmarks (left and right hand
tips, elbows, wrists, and thumbs, giving \(24\) spatial coordinates), split into
\(180\) training and \(180\) test sequences (\(30\) per gesture class in each
split). Because \(N_{\mathrm{train}} = 180\) (\(144\) per cross-validation fold)
while depth-two feature counts reach \(1{,}275\), both regularization tuning and
graph-based interaction restriction play a central role.

We construct paths from uncentered landmark displacements without temporal
resampling. Given landmark vectors \(h_r \in \R^{24}\) for frames
\(0 \leq r \leq 50\), we append a unit-duration clock channel \(r/50\) at channel
index \(0\) and prepend the origin anchor \(x_0 = 0 \in \R^{25}\):
\[
 x_0 = 0 \in \R^{25}, \qquad
 x_{r+1} = \Bigl(\frac{r}{50},\, h_r\Bigr) \in \R^{25}, \qquad
 z_r = x_{r+1} - x_r \quad (0 \leq r \leq 50).
\]
This produces \(R = 51\) increments: an initial-pose step \(z_0 = (0, h_0)\) with
zero clock advance, followed by \(50\) frame-to-frame displacement steps
\(z_r = (1/50,\, h_r - h_{r-1})\) for \(1 \leq r \leq 50\). The ordinary
first-level endpoint \(\sum_{r=0}^{50} z_r = (1, h_{50})\) therefore records the
final landmark configuration relative to the origin, while signed areas against
\(z_0\) retain sensitivity to the starting pose.

Applying the unmodified \(P = 2\) normalized Pascal lift~\eqref{eq:app-pascal}
with weights \(b_0(r) = 1\) and \(b_1(r) = u_r = r/50\) in copy-major order forms
\(50\)-channel lifted increments \(v_r = (z_r, u_r z_r) \in \R^{50}\), where
channels \(0, \ldots, 24\) are the ordinary increments and channels
\(25, \ldots, 49\) are their linearly position-weighted copies. All evaluated
maps use depth-two logsignatures over the \emph{entire} \(51\)-step sequence,
contrasting two undirected interaction graphs \(E\) on the \(C = 25\) input
channels via the dimension formula~\eqref{eq:graph-dimension}:
\begin{itemize}
\item \emph{Full interaction graph (\(E_{\mathrm{full}}\)):} Retains all
\(|E_{\mathrm{full}}| = \binom{25}{2} = 300\) unordered channel pairs. At
\(P = 1\), this gives \(25 + 300 = 325\) coordinates; at \(P = 2\),
\eqref{eq:graph-dimension} gives \(50 + 4(300) + 25\binom{2}{2} = 1{,}275 = 50 + \binom{50}{2}\)
coordinates.
\item \emph{Local interaction graph (\(E_{\mathrm{local}}\)):} Restricts
cross-channel edges to the \(24\) clock--spatial pairs \(\{0, c\}\)
(\(1 \leq c \leq 24\)) plus the \(8 \times \binom{3}{2} = 24\) within-landmark
\((x, y, z)\) coordinate pairs, giving \(|E_{\mathrm{local}}| = 48\) edges
(omitting cross-landmark spatial areas). Here ``local'' refers strictly to
kinematic channel locality, not temporal windowing. At \(P = 1\), this gives
\(25 + 48 = 73\) coordinates; at \(P = 2\),~\eqref{eq:graph-dimension} gives
\(50 + 4(48) + 25 = 267\) coordinates (\(50\) endpoints, \(192\) cross-channel
copy areas, and \(25\) same-channel separation areas \(J_{c_0 c_1}\)).
\end{itemize}

\subsection{Nested filtration of the nine feature maps}
Passing from an ordinary depth-two logsignature (\(P=1\)) to its \(P=2\) Pascal
counterpart adds three algebraically distinct groups of features:
\begin{enumerate}
\item the \(25\) position-weighted endpoints \(\sum_{r=0}^{50} u_r z_r\)
(\emph{moments});
\item \(73\) additional signed areas involving the weighted clock copy \(T_1\) or
clock--spatial copy pairings (bringing the total number of clock-related areas
from \(24\) to \(1 + 4 \times 24 = 97\), namely the clock self-area \(J_{T_0 T_1}\)
plus all four copy pairings across the \(24\) spatial channels); and
\item additional \emph{coded spatial--spatial areas} not involving the clock:
\(96\) in the local graph (\(24\) single-channel separation areas \(J_{c_0 c_1}\)
plus \(3 \times 24 = 72\) within-landmark cross-copy areas) and \(852\) in the
full graph (\(24\) single-channel separation areas plus \(3 \times \binom{24}{2} = 828\)
spatial cross-copy areas).
\end{enumerate}

To separate these contributions, we evaluate a nested four-step filtration of
coordinate sets within both the full and local graphs,
\[
 \text{Ordinary} \subset \text{Moments + ordinary}
 \subset \text{Coded-clock control} \subset \text{Pascal},
\]
together with the \(50\) coded endpoints alone (\emph{Moments only}), giving the
nine feature maps in Table~\ref{tab:natops-maps}. Nominal feature counts include
the constant clock coordinates \(s_{T_0} = 1\), \(s_{T_1} = 51/100\), and
\(J_{T_0 T_1}\) wherever retained. These map to zero under standardization. Nonconstant
dependent coordinates still contribute to the ridge Gram matrix. For each
spatial channel \(c\), for example,
\(J_{T_0 c_1} - 2J_{T_1 c_0} = s_{T_1}(s_{c_0} - s_{c_1})\)
(\nolinkurl{evidence/natops/verify_offset.py}).

\begin{table}[ht]
\centering
\begin{tabular}{@{}lrp{8.4cm}@{}}
\toprule
Feature map & Features & Retained endpoints and signed-area subsets \\
\midrule
Ordinary full & 325 & \(25\) ordinary endpoints \(+\) all \(300\) ordinary areas (\(24\) clock--spatial \(+\) \(276\) spatial). \\
Moments + ordinary full\(^*\) & 350 & \(50\) coded endpoints \(+\) all \(300\) ordinary areas. \\
Coded-clock control (full)\(^*\) & 423 & \(50\) coded endpoints \(+\) \(97\) coded-clock areas \(+\) \(276\) ordinary spatial areas. \\
Pascal full & 1,275 & \(50\) coded endpoints \(+\) all \(1{,}225\) coded areas (\(97\) clock \(+\) \(1{,}128\) coded spatial). \\
\midrule
Ordinary local & 73 & \(25\) ordinary endpoints \(+\) \(48\) ordinary local areas (\(24\) clock--spatial \(+\) \(24\) within-landmark). \\
Moments + ordinary local & 98 & \(50\) coded endpoints \(+\) \(48\) ordinary local areas. \\
Coded-clock control (local) & 171 & \(50\) coded endpoints \(+\) \(97\) coded-clock areas \(+\) \(24\) ordinary within-landmark areas. \\
Pascal local & 267 & \(50\) coded endpoints \(+\) \(217\) coded local areas (\(97\) clock \(+\) \(120\) coded local spatial). \\
\midrule
Moments only & 50 & \(50\) coded endpoints only (no signed areas). \\
\bottomrule
\end{tabular}
\caption{Coordinate filtration of the nine evaluated NATOPS depth-two
logsignature maps. \(^*\)Full-graph intermediate controls added after the initial
seven outcomes were observed, using the identical cross-validation rule.}
\label{tab:natops-maps}
\end{table}

\subsection{Shared cross-validation protocol, results, and paired analysis}
Every representation selects its ridge penalty on the exact same training
partition splits. Using five-fold stratified cross-validation on the \(180\)
training sequences (\(36\) validation sequences per fold, shuffled once with seed
\texttt{20261003}), we evaluate \(\alpha \in \{10^k : k = -6, \ldots, 6\}\) for a
one-versus-rest ridge classifier with intercept (Cholesky solver). Within each
fold, columnwise standardization (\texttt{StandardScaler}) is fitted on the \(144\)
fold-training sequences only. For each map, \(\alpha\) is chosen to maximize the
total number of correct validation predictions across all \(180\) training
sequences (equivalently mean five-fold validation accuracy), breaking exact ties
in favor of the smaller penalty \(\alpha\). The scaler and ridge classifier are
then refitted on all \(180\) training sequences at the selected \(\alpha\) and
evaluated once on the \(180\) test sequences. The seven unstarred maps and their
selected penalties were frozen prior to test evaluation; the two starred
full-graph controls were added afterward under the identical fold split and tie
rule. Table~\ref{tab:natops-results} reports both cross-validation and held-out
test results across all nine maps.

\begin{table}[ht]
\centering
\begin{tabular}{@{}lrrrrrr@{}}
\toprule
Feature map & Features & Selected \(\alpha\) & CV / 180 & CV (\%) & Test / 180 & Test (\%) \\
\midrule
Ordinary full & 325 & \(10\) & 166 & 92.22 & 163 & 90.56 \\
Moments + ordinary full\(^*\) & 350 & \(10\) & 168 & 93.33 & 166 & 92.22 \\
Coded-clock control (full)\(^*\) & 423 & \(100\) & 166 & 92.22 & 163 & 90.56 \\
Pascal full & 1,275 & \(10\) & 166 & 92.22 & 173 & 96.11 \\
\midrule
Ordinary local & 73 & \(10\) & 167 & 92.78 & 167 & 92.78 \\
Moments + ordinary local & 98 & \(1\) & 167 & 92.78 & 165 & 91.67 \\
Coded-clock control (local) & 171 & \(10\) & 169 & 93.89 & 168 & 93.33 \\
Pascal local & 267 & \(10\) & 169 & 93.89 & 165 & 91.67 \\
\midrule
Moments only & 50 & \(0.1\) & 143 & 79.44 & 142 & 78.89 \\
\bottomrule
\end{tabular}
\caption{NATOPS cross-validation selections and held-out test performance across
all nine maps (\cite[Table~2]{mainpaper}). \(^*\)Controls added after the initial
seven test outcomes.}
\label{tab:natops-results}
\end{table}

Three empirical findings emerge from Table~\ref{tab:natops-results}:
\begin{enumerate}
\item \emph{Performance of the endpoints-only control.} The \(50\) ordinary and
position-weighted endpoints (\emph{Moments only}) achieve \(79.44\%\) CV and
\(78.89\%\) test accuracy (\(142/180\)), below the maps that include signed areas
on this split.
\item \emph{Divergence between the full and local interaction graphs:} On the
full interaction graph, Pascal coding increases test accuracy from \(90.56\%\)
(\(163/180\) for \emph{Ordinary full}, \(\alpha = 10\)) to \(96.11\%\)
(\(173/180\) for \emph{Pascal full}, \(\alpha = 10\)). Neither adding weighted
endpoints alone (\(166/180\), \(92.22\%\)) nor adding all \(97\) coded-clock
areas (\(163/180\), \(90.56\%\), selecting \(\alpha = 100\)) reproduces that
gain; at the sequence level, \emph{Pascal full} corrects \(12\) errors of
\emph{Ordinary full} while introducing \(2\) (\(161\) both correct, \(5\) both
wrong; descriptive exact two-sided McNemar \(p = 53/4096 \approx 0.0129\)), and
corrects \(11\) errors of \emph{Coded-clock control (full)} while introducing
\(1\). On the restricted local graph, however, \emph{Pascal local} scores
\(165/180\) (\(91.67\%\)) against \(167/180\) (\(92.78\%\)) for \emph{Ordinary
local} and \(168/180\) (\(93.33\%\)) for \emph{Coded-clock control (local)}. Four
pre-specified paired comparisons of \emph{Pascal local} against \emph{Ordinary
local}, \emph{Ordinary full}, \emph{Moments + ordinary local}, and
\emph{Coded-clock control (local)} give unadjusted two-sided exact McNemar
\(p\)-values between \(0.55\) and \(1.00\), and \(1.00\) after Holm adjustment.
\item \emph{Cross-validation ranking and regularization confounding:} Displaying
CV and test scores side by side highlights an important small-sample caveat. On
the \(180\) training sequences, five-fold CV assigns \(92.22\%\) (\(166/180\)) to
\emph{Pascal full}, slightly \emph{below} the local maps (\(92.78\%\)--\(93.89\%\));
thus training-only model selection across graphs would choose a local map rather
than \emph{Pascal full}. Furthermore, passing from \emph{Coded-clock control
(full)} (\(423\) features, \(\alpha = 100\)) to \emph{Pascal full} (\(1{,}275\)
features, \(\alpha = 10\)) simultaneously introduces both within-landmark and
cross-landmark coded spatial areas, triples the ambient dimension, and shifts the
selected penalty by an order of magnitude. The contrast therefore points to
informative coded spatial--spatial interactions in the full map, without
isolating cross-landmark edges from feature width and ridge geometry.
\end{enumerate}
\FloatBarrier

\FloatBarrier\Needspace{10\baselineskip}
\section{Cross-subject transfer, split protocols, and computational cost}
\label{app:additional}

This section brings together two additional studies that complement the UCI HAR
and NATOPS experiments in~\cite{mainpaper}: a nested cross-participant study on
a 15-sensor body-worn dataset (DSADS), and an examination of subject-mixed
window splitting and single-thread CPU pipeline runtime on UCI HAR.

\subsection{Held-out participants on DSADS}\label{app:transfer}
Does combining QR-Pascal temporal and spectral logsignatures generalize to unseen
participants on a multi-sensor wearable benchmark without smartphone-specific
gravity decomposition or hand-tuned lag features? We evaluate a generic two-path
per-sensor recipe on the Daily and Sports Activities Dataset (DSADS)~\cite{dsads},
which records \(19\) activities across \(8\) participants in five-second segments
of \(R = 125\) samples at \(25\,\mathrm{Hz}\) (\(60\) segments per participant and
activity, totaling \(9{,}120\) segments, or \(1{,}140\) per participant). Five
body-worn units (torso, right arm, left arm, right leg, and left leg) each carry
a three-axis accelerometer, gyroscope, and magnetometer, yielding \(15\)
three-axis sensors (\(45\) scalar channels).

\paragraph{Generic two-path per-sensor representation.}
Within each outer training split, every scalar channel is standardized using its
training mean and training population standard deviation (floored at \(10^{-8}\)),
and passed pointwise through the inverse hyperbolic sine. Writing \(a_r\) for a
standardized reading, this gives
\[
 y_r = \arsinh(a_r) = \log\bigl(a_r + \sqrt{a_r^2 + 1}\bigr)
 \qquad (0 \leq r < 125),
\]
which acts linearly near zero while logarithmically compressing heavy-tailed
impact spikes without saturating to a finite bound. Grouping the \(45\) channels
into their \(15\) physical three-axis sensors \(y_r^{(g)} \in \R^3\)
(\(1 \leq g \leq 15\)), we construct two \(9\)-channel lifted paths per sensor
using the \(P = 3\) orthonormal QR-Pascal weights \(q_p\) of Section~\ref{app:qr}:
\begin{enumerate}
\item \emph{Temporal sensor path:} Treating scaled readings
\(y_r^{(g)} / 125 \in \R^3\) as increments over \(r = 0, \ldots, 124\) and
lifting by \((q_0(r), q_1(r), q_2(r))\) produces \(9\) channels and
\(9 + \binom{9}{2} = 45\) depth-two logsignature coordinates, preserving static
sensor levels alongside within-sensor cross-axis temporal areas.
\item \emph{Spectral sensor path:} For each of the three transformed axes of
sensor \(g\), we compute Parseval-scaled one-sided Fourier bin powers over
\(K = 63\) bins (\(k = 0, \ldots, 62\)) and normalize each axis by its sum of bin
powers (setting all shares to zero if total axis power is at most \(10^{-20}\)).
Lifting the three frequency-ordered share vectors jointly by
\((q_0(k), q_1(k), q_2(k))\) produces \(9\) channels and \(45\) nominal depth-two
logsignature coordinates (\(42\) after dropping the three constant-copy endpoints,
which equal \(1/\sqrt{63}\) when all three axes have power above \(10^{-20}\)).
\end{enumerate}
Because spectral normalization removes amplitude scale, each sensor \(g\) also
contributes one log-energy summary computed before standardization,
\[
 \log\Bigl(1 + \frac{E_g}{E_{g,\mathrm{train}}}\Bigr), \qquad
 E_g = \frac{1}{125}\sum_{r=0}^{124} \|x_r^{(g)}\|^2,
\]
where \(x_r^{(g)} \in \R^3\) is the uncentered raw triad and \(E_{g,\mathrm{train}}\)
is the training mean of \(E_g\) (with \(\sqrt{E_{g,\mathrm{train}}}\) floored at
\(10^{-8}\)). Concatenating across all \(15\) sensors gives
\(15 \times (45 + 42 + 1) = 1{,}320\) retained features (\(1{,}365\) nominal).

\paragraph{Nested participant-disjoint evaluation and ablations.}
Each of the eight participants serves in turn as the held-out outer test cohort
(\(1{,}140\) segments), while the remaining seven participants (\(7{,}980\)
segments) form the outer training set. Within each seven-person training cohort,
leave-one-participant-out cross-validation (fitting on six people and validating
on the seventh) selects the mean-squared-error ridge penalty
\(\lambda \in \{10^{-8}, 10^{-7}, \ldots, 10^2, \infty\}\) (applied after
training-only variance masking, column standardization, and \(1/\sqrt{d_{\mathrm{ret}}}\)
width scaling) jointly with a binary indicator of whether query features are
clipped to their training column ranges, before refitting on all seven training
participants. Table~\ref{tab:dsads} compares the full \(P=3\) representation
against two matched ablations that keep the identical \(\arsinh\) preprocessing,
energy summaries, and nested tuning protocol: dropping all signed areas (retaining
only coded endpoints and energy) and setting \(P = 1\) (retaining within-sensor
cross-axis areas with only the constant weight \(q_0\)).

\begin{table}[ht]
\centering
\begin{tabular}{@{}lrrrrr@{}}
\toprule
DSADS representation (\(\arsinh\) + energy) & Features & Correct & Errors & Accuracy (\%) & SD (pp) \\
\midrule
Full \(P=3\) temporal + spectral logsignature & 1,320 & 8,466 & 654 & 92.83 & 1.20 \\
Endpoints only (drop signed areas, \(P=3\)) & 240 & 8,122 & 998 & 89.06 & 3.75 \\
Constant weight only (\(P=1\), keep areas) & 150 & 8,131 & 989 & 89.16 & 4.47 \\
\bottomrule
\end{tabular}
\caption{Nested eight-participant leave-one-subject-out results on DSADS across
all \(8 \times 1{,}140 = 9{,}120\) test segments. ``Features'' counts retained
columns; ``SD (pp)'' is the sample standard deviation of test accuracy (in
percentage points) across the eight held-out participants.}
\label{tab:dsads}
\end{table}

Across all \(9{,}120\) held-out segments, the combined \(P = 3\) logsignature
reaches \(92.83\%\) accuracy (\(654\) errors, between-participant SD \(1.20\)
percentage points). Omitting signed areas lowers accuracy by \(3.77\) percentage
points to \(89.06\%\) (\(998\) errors), while omitting the linear and quadratic
position weights (\(P = 1\)) lowers accuracy by \(3.67\) percentage points to
\(89.16\%\) (\(989\) errors), indicating that both ordered cross-channel areas and
nonconstant position weights contribute materially on average across wearers. As
documented in the archived transform study (\nolinkurl{evidence/har/transfer/transforms/}),
this cohort had previously informed feature design, and selecting the
representation family itself across a wider \(68\)-family candidate menu via the
seven inner validation participants yields \(91.33\%\) pooled accuracy (\(791\)
errors, SD \(3.93\) pp), reflecting ranking variability when choosing among many
feature families from seven validation subjects.

\subsection{Subject-mixed UCI HAR, SpikeDCL protocol, and pipeline runtime}
\label{app:compute}
Published activity-recognition benchmarks do not all evaluate on held-out
participants, and comparing accuracies or runtimes across papers requires
tracking both the data partition and the stages included in timing measurements.

\paragraph{Subject-mixed random-window split versus held-out subjects.}
Whereas Section~\ref{app:har} holds out nine of the \(30\) UCI HAR subjects,
SpikeDCL---a spatio-temporal spiking neural network for activity
recognition~\cite{spikedcl}---evaluates on a subject-mixed random-window
partition. Its loader concatenates all \(7{,}352 + 2{,}947 = 10{,}299\) windows
ordered by subject, draws \(20\%\) (\(2{,}060\) windows) at random without
stratification (seed \(0\)) as the test set, and draws \(20\%\) of the remainder
(\(1{,}648\) windows, seed \(0\)) for validation, leaving \(6{,}591\) training
windows. Two consequences distinguish this protocol from subject-disjoint
transfer: all \(30\) participants appear in every split, and because raw UCI HAR
windows of length \(128\) advance with a \(50\%\) stride (\(64\) samples),
adjacent windows placed in training and test partitions share half of their
underlying sensor readings.

Evaluating the unchanged \emph{compact368} representation on this exact
\(6{,}591 / 2{,}060\) partition---fitting observation calibration, the variance
mask, and feature standardization strictly on the \(6{,}591\) training windows,
with fixed \(\alpha = 10\) and inverse-frequency sample weights
\(6591 / (6 n_c)\) for class \(c\) matching the SpikeDCL loader's class
balancing---achieves \(98.16\%\) test accuracy (\(38\) errors out of \(2{,}060\)
windows). SpikeDCL reports \(98.86 \pm 0.28\%\) across five neural training
seeds~\cite[Table~2]{spikedcl}. These subject-mixed numbers belong to a
qualitatively easier generalization regime than the held-out-subject scores in
Section~\ref{app:har}, and sharing a split does not equate model families,
architectural inductive biases, or hyperparameter tuning budgets.

\paragraph{Single-thread CPU runtime of the \emph{compact368} pipeline.}
Because lifting \(C\) channels by \(P\) copies expands width to \(PC\), low
signature truncation depth (\(m = 2\)) does not automatically guarantee fast
extraction. Nevertheless, depth-two evaluation requires only a single linear pass
of prefix sums and outer products per window, and fitting~\eqref{eq:har-ridge} is
a single \(368 \times 368\) linear solve. To document concrete wall-clock costs
on the \(6{,}591\)-train / \(2{,}060\)-test SpikeDCL partition, we timed every
stage of \emph{compact368} on a single CPU thread (Apple M5 Pro, Python~3.13,
PyTorch~2.13, NumPy~2.5, scikit-learn~1.9; one full warm-up run followed by five
measured repeats, excluding disk I/O). An automated audit verified zero
numerical discrepancy between the timed pipeline and the archived evaluation on
all training features, test features, and test predictions.

\begin{table}[ht]
\centering
\begin{tabular}{@{}lrr@{}}
\toprule
\emph{compact368} pipeline stage (single CPU thread, milliseconds) & Median & Min--Max \\
\midrule
\multicolumn{3}{@{}l}{\emph{Training pipeline (\(6{,}591\) windows)}} \\
Observation calibration and preprocessing & 174.47 & 162.10--179.85 \\
Feature extraction (raw, spectral, and geometric branches) & 1,304.49 & 1,300.74--1,313.52 \\
Variance mask and feature standardization & 7.88 & 7.79--8.51 \\
Closed-form weighted ridge solve (\(\alpha = 10\)) & 4.49 & 4.43--4.75 \\
Total training time & 1,491.17 & 1,488.76--1,493.80 \\
\midrule
\multicolumn{3}{@{}l}{\emph{Test inference (\(2{,}060\) windows)}} \\
Observation preprocessing & 31.42 & 29.39--42.03 \\
Feature extraction & 422.19 & 417.19--446.65 \\
Frozen variance mask and feature standardization & 0.85 & 0.84--1.29 \\
Linear score evaluation & 0.09 & 0.09--0.10 \\
Total inference time & 456.94 & 447.52--479.46 \\
\bottomrule
\end{tabular}
\caption{Single-thread CPU wall-clock times (in ms across five repeats after one
warm-up, excluding disk archive loading of \(1{,}165.43\,\mathrm{ms}\)) for
\emph{compact368} on the \(6{,}591 / 2{,}060\)-window HAR partition.}
\label{tab:compute}
\end{table}

Table~\ref{tab:compute} shows that end-to-end training over \(6{,}591\) windows
completes in a median of \(1.49\,\mathrm{s}\)---dominated by feature extraction
(\(1.30\,\mathrm{s}\)), with the ridge solve itself requiring \(4.49\,\mathrm{ms}\)---and
batch inference over all \(2{,}060\) test windows takes \(0.46\,\mathrm{s}\)
(\(0.22\,\mathrm{ms}\) amortized per window). By contrast, SpikeDCL trains a
spiking network for \(60\) gradient epochs across three candidate learning
rates~\cite[Sec.~3.1]{spikedcl}. Reporting these batch CPU timings documents the
practical cost of the coded signature pipeline without claiming a hardware-matched
speedup or single-window embedded latency comparison against~\cite{spikedcl}.
\FloatBarrier


\clearpage
\phantomsection
\section*{Reproducing the results}\label{app:reproducibility}
\addcontentsline{toc}{section}{Reproducing the results}

The repository
\href{https://github.com/Froskekongen/pascal_coded_application}{\nolinkurl{github.com/Froskekongen/pascal_coded_application}}
organizes verification and empirical evidence into two complementary tiers:
\begin{enumerate}
\item \emph{Clean-room executable reproduction (\nolinkurl{reproduction/}):}
Documented step by step in
\href{https://github.com/Froskekongen/pascal_coded_application/blob/main/reproduction/README.md}{\nolinkurl{reproduction/README.md}}
and locked via \nolinkurl{pyproject.toml} and \nolinkurl{uv.lock} using
\texttt{tensorsig}~\cite{tensorsig}. Running \texttt{uv sync --locked},
\texttt{uv run python -m reproduction.checks},
\texttt{uv run python -m reproduction.uci\_har --download}, and
\texttt{uv run python -m reproduction.natops --download} from the repository root
runs offline feature-identity checks, verifies SHA-256 dataset checksums, refits
the three primary official-split UCI HAR models (\emph{compact368},
\emph{geometry57}, and the \(561\)-feature baseline) and all nine NATOPS models
from scratch, and compares confusion matrices, test prediction vectors, and all
\(585\) NATOPS cross-validation fold counts against frozen expected outputs in
the \nolinkurl{reproduction/reference/} directory.
\item \emph{Mathematical checks and archived study records}
(\nolinkurl{paper/source/evidence/}). Standard-library Python scripts check
identities from Part~I without datasets or fitted models, using rational
arithmetic for the Pascal and offset checks and binary64 for the spectral
checks. Packaged CSV/JSON records, source snapshots, and audit scripts preserve the
\(187\)-configuration HAR conventional-signature comparison, the nested
eight-participant DSADS transfer study, and the single-thread CPU timing
benchmark.
\end{enumerate}

Table~\ref{tab:repro-index} maps every figure, theorem, and empirical table across
the main paper~\cite{mainpaper} and this companion to its corresponding
code and evidence files.

\begin{table}[ht]
\centering
\small
\begin{tabular}{@{}>{\raggedright\arraybackslash}p{4.3cm}>{\raggedright\arraybackslash}p{4.6cm}>{\raggedright\arraybackslash}p{6.1cm}@{}}
\toprule
Paper / companion item & Executable script or source & Archived evidence / verification record \\
\midrule
Main Fig.~1; Companion Fig.~\ref{fig:comp-movement} \& Ex.~\ref{ex:two-rises} (two-rise schedules, \(1/2\) vs.\ \(1/6\))
  & \nolinkurl{evidence/natops/make_movement_figure.py}; \nolinkurl{reproduction/checks.py}
  & \nolinkurl{evidence/natops/verify_offset.py}; \nolinkurl{evidence/natops/verification_offset.json} \\[4pt]
Companion Sec.~\ref{app:offsets}, Props.~\ref{prop:qr-uniqueness}--\ref{prop:nilpotent} (QR, shifts, twisted Chen)
  & \nolinkurl{evidence/natops/verify_pascal.py}; \nolinkurl{evidence/natops/verify_spectral.py}
  & \nolinkurl{evidence/natops/verification_pascal.json}; \nolinkurl{evidence/natops/verification_spectral.json} \\[4pt]
Companion Sec.~\ref{app:depth-and-ridge}, Ex.~\ref{ex:depth-three-collision}, Thm.~\ref{prop:staircase-projection} (depth-3 collision \& staircase projection)
  & \nolinkurl{evidence/natops/verify_offset.py}; \nolinkurl{evidence/har/historical_selected_sources/ordinary_spectral.py}
  & \nolinkurl{evidence/natops/verification_offset.json}; \nolinkurl{evidence/har/consolidated_results/results.csv} \\[4pt]
Main Table~1 (\emph{compact368}, supplied 561); Companion Tables~\ref{tab:compact-blocks}--\ref{tab:har-standalone} (\emph{geometry57})
  & \nolinkurl{reproduction/uci_har.py}; \nolinkurl{evidence/har/retained_features.py}; \nolinkurl{evidence/har/geometry57.py}
  & \nolinkurl{reproduction/reference/har.json}; \nolinkurl{evidence/har/retained_results/}; \nolinkurl{evidence/har/geometry57_results/} \\[4pt]
Main Table~1 (controls); Companion Tables~\ref{tab:har-standalone}--\ref{tab:ordinary-spectral} (standalone, deletions, raw \& staircase controls)
  & \nolinkurl{evidence/har/historical_selected_sources/ordinary_controls.py}; \nolinkurl{ordinary_spectral.py}
  & \nolinkurl{evidence/har/pool_results/results.csv}; \nolinkurl{evidence/har/consolidated_results/results.csv}; \nolinkurl{saved_metrics_verification.json} \\[4pt]
Main Table~2; Companion Sec.~\ref{app:natops}, Tables~\ref{tab:natops-maps}--\ref{tab:natops-results} (nine NATOPS maps, CV \& McNemar)
  & \nolinkurl{reproduction/natops.py}; \nolinkurl{evidence/natops/offset_learning/reproduce.py}
  & \nolinkurl{reproduction/reference/natops.json}; \nolinkurl{evidence/natops/offset_learning/results.json}; \nolinkurl{diagnostic_full_results.json} \\[4pt]
Companion Sec.~\ref{app:transfer}, Table~\ref{tab:dsads} (nested 8-participant DSADS transfer \& ablations)
  & \nolinkurl{evidence/har/transfer/transforms/source_snapshot.zip} (\nolinkurl{streams.py}, \nolinkurl{study.py})
  & \nolinkurl{evidence/har/transfer/transforms/summary/dataset_overview.csv}; \nolinkurl{dsads_model_distribution.csv}; \nolinkurl{results/verification.json} \\[4pt]
Companion Sec.~\ref{app:compute}, Table~\ref{tab:compute} (SpikeDCL random-window split \& CPU runtime)
  & \nolinkurl{evidence/timing/benchmark_pipeline.py}; \nolinkurl{evidence/timing/frozen_accuracy/sources/}
  & \nolinkurl{evidence/timing/summary.csv}; \nolinkurl{evidence/timing/timings.csv}; \nolinkurl{evidence/timing/frozen_accuracy/run.json} \\
\bottomrule
\end{tabular}
\caption{Index connecting figures, mathematical checks, and experimental tables in
the main paper~\cite{mainpaper} and this companion to executable scripts and
archived verification records (paths under \nolinkurl{evidence/} are relative to
\nolinkurl{paper/source/}).}
\label{tab:repro-index}
\end{table}
\FloatBarrier

\input{common/ai_disclosure.tex}
\Needspace{8\baselineskip}
\begin{thebibliography}{99}
\phantomsection
\addcontentsline{toc}{section}{References}
\bibitem{mainpaper} E. Aune and F. N. Harang. Building activity representations with Pascal-coded signatures. Preprint, 2026.
\bibitem{chen} K.-T. Chen. Integration of paths, geometric invariants and a generalized Baker--Hausdorff formula. \href{https://doi.org/10.2307/1969671}{\emph{Annals of Mathematics}}, 65:163--178, 1957.
\bibitem{yang} W. Yang, T. Lyons, H. Ni, C. Schmid, and L. Jin. Developing the path signature methodology and its application to landmark-based human action recognition. In \href{https://doi.org/10.1007/978-3-030-98519-6_18}{\emph{Stochastic Analysis, Filtering, and Stochastic Optimization}}, pp.~431--464, Springer, 2022.
\bibitem{generalised} J. Morrill, A. Fermanian, P. Kidger, and T. Lyons. A generalised signature method for multivariate time series feature extraction. \href{https://arxiv.org/abs/2006.00873v2}{\emph{arXiv:2006.00873v2}}, 2021.
\bibitem{fruits} J. Diehl and R. Krieg. FRUITS: feature extraction using iterated sums for time series classification. \href{https://doi.org/10.1007/s10618-024-01068-1}{\emph{Data Mining and Knowledge Discovery}}, 38:4122--4156, 2024.
\bibitem{volterra} P. P. Hager, F. N. Harang, L. Pelizzari, and S. Tindel. The Volterra signature. \href{https://arxiv.org/abs/2603.04525v1}{\emph{arXiv:2603.04525v1}}, 2026.
\bibitem{ews} A. Bloch, B. Walker, J. Mouterde, S. Morley, S. N. Cohen, and T. Lyons. Extending SSMs with the exponentially weighted signature. \href{https://arxiv.org/abs/2603.19198v3}{\emph{arXiv:2603.19198v3}}, 2026.
\bibitem{gini} G. A. Koshevoy and K. Mosler. Multivariate Gini indices. \href{https://wisostat.uni-koeln.de/sites/dep_econometrics/pdf_publikationen/1997_KoshevoyM97a_Gini.pdf}{\emph{Journal of Multivariate Analysis}}, 60:252--276, 1997.
\bibitem{polynomialoneforms} H. Boedihardjo and X. Geng. The uniqueness of signature problem in the non-Markov setting. \href{https://arxiv.org/abs/1401.6165v2}{\emph{arXiv:1401.6165v2}}, 2014.
\bibitem{discreteqr} L. Reichel. Fast QR decomposition of Vandermonde-like matrices and polynomial least squares approximation. \href{https://doi.org/10.1137/0612041}{\emph{SIAM Journal on Matrix Analysis and Applications}}, 12(3):552--564, 1991.
\bibitem{ridge} C. Saunders, A. Gammerman, and V. Vovk. Ridge regression learning algorithm in dual variables. In \href{https://pure.royalholloway.ac.uk/en/publications/ridge-regression-learning-algorithm-in-dual-variables/}{\emph{Proceedings of the Fifteenth International Conference on Machine Learning}}, pp.~515--521, Morgan Kaufmann, 1998.
\bibitem{sigadapt} N. Drobac, M. Br\'eg\`ere, J. de Vilmarest, and O. Wintenberger. Encoding time series dynamics with sliding-window signatures for online forecasting. \href{https://arxiv.org/abs/2610.35367v1}{\emph{arXiv:2610.35367v1}}, 2026.
\bibitem{har} D. Anguita, A. Ghio, L. Oneto, X. Parra, and J. L. Reyes-Ortiz. A public domain dataset for human activity recognition using smartphones. \emph{ESANN}, pp.~437--442, 2013.
\bibitem{natopsoriginal} N. Ghouaiel, P.-F. Marteau, and M. Dupont. Continuous pattern detection and recognition in stream---a benchmark for online gesture recognition. \href{https://doi.org/10.1504/IJAPR.2017.085315}{\emph{International Journal of Applied Pattern Recognition}}, 4(2):146--160, 2017.
\bibitem{natops} TSML Archive Maintainers. NATOPS dataset. UEA time series classification archive. \href{https://doi.org/10.5281/zenodo.11206248}{Zenodo record 11206248}, 2018.
\bibitem{dsads} B. Barshan and K. Altun. \emph{Daily and Sports Activities} [Dataset]. UCI Machine Learning Repository, 2010. \href{https://doi.org/10.24432/C5C59F}{doi:10.24432/C5C59F}.
\bibitem{spikedcl} Y. Li, R. Yin, Y. Kim, and P. Panda. Efficient human activity recognition with spatio-temporal spiking neural networks. \href{https://doi.org/10.3389/fnins.2023.1233037}{\emph{Frontiers in Neuroscience}}, 17:1233037, 2023.
\bibitem{tensorsig} E. Aune. \emph{tensorsig} (Version 0.1.0) [Computer software]. Zenodo, 2026. \href{https://doi.org/10.5281/zenodo.23886748}{doi:10.5281/zenodo.23886748}.
\end{thebibliography}

\end{document}
